\subsection{伴随表示}

设 \(G\) 是一个李群。考虑解析左运算律

\[
(g, g ^ {\prime}) \mapsto g g ^ {\prime} g ^ {- 1} = (\operatorname{Int} g) g ^ {\prime}
\]

从 G 到 G。根据第3号，此运算律定义了一个从 \(\mathcal{T}^{(\infty)}(\mathrm{G})\times \mathcal{T}^{(\infty)}(\mathrm{G})\) 到 \(\mathcal{T}^{(\infty)}(\mathrm{G})\) 的双线性映射，本号中记为 \(\top\)。根据第3号命题13，

\[
(t * t ^ {\prime}) \top t ^ {\prime \prime} = t \top (t ^ {\prime} \top t ^ {\prime \prime})\tag{22}
\]

对任意 \(t, t', t''\)，它们属于 \(\mathcal{T}^{(\infty)}(\mathrm{G})\)。根据第3号命题14 (i)，

\[
\varepsilon_ {g} \top t = (\operatorname{Int} g) _ {*} t\tag{23}
\]

对任意 \(g \in G\) 及 \(t \in \mathcal{T}^{(\infty)}(G)\)。特别地，映射 \(t \mapsto \varepsilon_{g} \top t\) 是从 \(\mathcal{T}^{(\infty)}(G)\) 到 \(\mathcal{T}^{(\infty)}(G)\) 的双代数 \(\mathcal{T}^{(\infty)}(G)\) 的自同构。它在 \(U(G)\)、\(U_{s}(G)\)、\(L(G)\) 上的限制分别记为 \(\mathrm{Ad}_{\mathrm{U(G)}}(g)\)、\(\mathrm{Ad}_{\mathrm{U_s(G)}}(g)\)、\(\mathrm{Ad}_{\mathrm{L(G)}}(g)\)。不致混淆时，我们常用 \(\mathrm{Ad}(g)\) 代替 \(\mathrm{Ad}_{\mathrm{L(G)}}(g)\)。由 (23)，\(\mathrm{Ad}(g)\) 是在 \(e\) 处的 \(\operatorname{Int}(g)\) 的切映射。它是可赋范

李代数 \(\mathbf{L}(\mathbf{G})\) 的自同构。当 \(\mathbf{K}\) 的特征为 0 时，\(\operatorname{Ad}_{\mathbf{U}(\mathbf{G})}(g)\) 是 \(\mathbf{U}(\mathbf{G})\) 的唯一自同构，它延拓 \(\operatorname{Ad}(g)\)。

若 \(\phi\) 是从李群 G 到李群 H 的态射，则

\[
\phi_ {*} (t \top t ^ {\prime}) = \phi_ {*} (t) \top \phi_ {*} (t ^ {\prime})\tag{24}
\]

对任意 \(t, t'\)，它们属于 \(\mathcal{T}^{(\infty)}(\mathrm{G})\)；这由第3号命题15推出。

命题 43. 设 \(t, u\) 属于 \(\mathcal{T}^{(\infty)}(\mathrm{G})\)。设 \(\sum_{i=1}^{n} t_i \otimes t'_i\) 是 \(t\) 在余乘法下的像。则

\[
t \top u = \sum_ {i = 1} ^ {n} t _ {i} * u * t _ {i} ^ {\prime \vee}.
\]

按定义，\(t \top u\) 是 \(t \otimes u\) 在映射 \((g, g') \mapsto gg'g^{-1}\) 下的像，该映射从 \(G \times G\) 到 \(G\)。现在，这个映射由下列映射复合得到：

\[
\begin{array}{l l} \alpha \colon (g, g ^ {\prime}) \mapsto (g, g, g ^ {\prime}) & \text {of G \times G into G \times G \times G} \\ \beta \colon (g, g ^ {\prime}, g ^ {\prime \prime}) \mapsto (g, g ^ {\prime - 1}, g ^ {\prime \prime}) & \text {of G \times G \times G into G \times G \times G} \\ \gamma \colon (g, g ^ {\prime}, g ^ {\prime \prime}) \mapsto g g ^ {\prime \prime} g ^ {\prime} & \text {of G \times G \times G into G}. \end{array}
\]

另一方面：

\[
\begin{array}{c} \alpha_ {*} (t \otimes u) = \sum_ {i = 1} ^ {n} (t _ {i} \otimes t _ {i} ^ {\prime}) \otimes u = \sum_ {i = 1} ^ {n} t _ {i} \otimes t _ {i} ^ {\prime} \otimes u \\ \beta_ {*} (\sum_ {i = 1} ^ {n} t _ {i} \otimes t _ {i} ^ {\prime} \otimes u) = \sum_ {i = 1} ^ {n} t _ {i} \otimes t _ {i} ^ {\prime \vee} \otimes u \\ \gamma_ {*} (\sum_ {i = 1} ^ {n} t _ {i} \otimes t _ {i} ^ {\prime \vee} \otimes u) = \sum_ {i = 1} ^ {n} t _ {i} * u * t _ {i} ^ {\prime \vee}. \end{array}
\]

推论 1. 设 \(u \in \mathrm{L}(\mathbf{G})\)，且 \(u' \in \mathcal{T}^{(\infty)}(\mathbf{G})\)。则 \(u \top u' = u * u' - u' * u\)。

u 在余乘法下的像是 \(u \otimes \varepsilon_{e} + \varepsilon_{e} \otimes u\)，所以

\[
u \top u ^ {\prime} = u * u ^ {\prime} * \varepsilon_ {e} + \varepsilon_ {e} * u ^ {\prime} * u ^ {\vee} = u * u ^ {\prime} - u ^ {\prime} * u.
\]

推论 2. 设 \(t \in \mathcal{T}^{(\infty)}(\mathrm{G})\) 且 \(g \in \mathrm{G}\)。则 \(\varepsilon_g \top t = \varepsilon_g * t * \varepsilon_g^{-1}\)。若 \(t \in \mathrm{L}(\mathrm{G})\)，则 \(\varepsilon_g \top t = gtg^{-1}\)（后一个乘积在群 \(\mathrm{T}(\mathrm{G})\) 中计算）。

\(\varepsilon_{g}\) 在余乘法下的像是 \(\varepsilon_{g} \otimes \varepsilon_{g}\)。

推论 3. 设 \(a \in \mathrm{L}(\mathrm{G})\)。由 \(a\) 及左运算 \(g \mapsto \operatorname{Int} g\) 定义的 \(\mathrm{G}\) 在 \(\mathrm{G}\) 上的向量场是场 \(\mathbf{R}_a - \mathbf{L}_a\)。

这个场在 \(g\) 处的值为

\[
\begin{array}{r l} a \top \varepsilon_ {g} = a * \varepsilon_ {g} - \varepsilon_ {g} * a & (\text { Corollary   1 }) \\ = (R _ {a}) _ {g} - (L _ {a}) _ {g} & (\text { Definition   5 }). \end{array}
\]

对任意 \(g \in G\) 及 \(t \in \mathrm{L}(\mathrm{G})\)，

\[
(\mathrm{Ad} g) (t) = \varepsilon_ {g} \top t = \varepsilon_ {g} * t * \varepsilon_ {g ^ {- 1}} = g t g ^ {- 1}.\tag{25}
\]

由于 Ad \(g = \mathrm{T}_{e}(\mathrm{Int} g)\)，第11号命题42证明了 Ad 是 G 在可赋范空间 \(\mathrm{L}(\mathrm{G})\) 上的解析线性表示。

定义 8. G 在 L(G) 上的表示 Ad 称为 G 的伴随表示。

命题 44. 对任意 \(a \in \mathbf{L}(\mathbf{G})\)，

\[
(\mathrm{L} (\mathrm{Ad})) (a) = \mathrm{ad} _ {\mathrm{L} (\mathrm{G})} a.
\]

设 \(b \in \mathrm{L}(\mathrm{G})\)。根据第11号命题42 (ii)及命题43的推论3，

\[
(\mathrm{L} (\mathrm{Ad})) (a). b = - [ \mathrm{R} _ {a} - \mathrm{L} _ {a}, \mathrm{L} _ {b} ] (e).
\]

现在 \(R_{a} \circ L_{b} = L_{b} \circ R_{a}\)（第6号命题23 (ii)），因而 \([R_{a}, L_{b}] = 0\)；再利用命题23 (ii)，

\[
(\mathrm{L} (\mathrm{Ad})) (a). b = [ \mathrm{L} _ {a}, \mathrm{L} _ {b} ] (e) = \mathrm{L} _ {[ a, b ]} (e) = [ a, b ] = (\mathrm{ad} _ {\mathrm{L} (\mathrm{G})} a) b.
\]

命题 45. 假设 G 是有限维的，且 K 的特征为 0。设 s 是满足 \(\geqslant0\) 的整数。则映射 \(\pi: g \mapsto \operatorname{Ad}_{\mathrm{U}_{s}(\mathbf{G})}(g)\) 是 G 在 \(\mathrm{U}_{s}(\mathrm{G})\) 上的解析线性表示，并且 \(\mathrm{L}(\pi)a = \operatorname{ad}_{\mathrm{U}_{s}(\mathbf{G})}a\) 对任意 \(a \in \mathrm{L}(\mathrm{G})\) 成立。

线性表示 \(\pi\) 是 \(\bigoplus_{r=0}^{\infty} \mathrm{T}^{r}(\mathrm{Ad})\) 的一个商，因而是解析的。对于 \(a \in \mathrm{L}(\mathrm{G})\) 以及 \(x_{1}, x_{2}, \ldots, x_{s}\)，它们属于 \(\mathrm{L}(\mathrm{G})\)，

\[
\begin{array}{l l} (\mathrm{L} (\pi) a) (x _ {1} x _ {2} \dots x _ {s}) = \sum_ {i = 1} ^ {s} x _ {1} \dots (\mathrm{L} (\mathrm{Ad}) a. x _ {i}) \dots x _ {s} & \text {(Proposition 41)} \\ = \sum_ {i = 1} ^ {s} x _ {1} \dots ([ a, x _ {i} ]) \dots x _ {s} & \text {(Proposition 44)} \\ = (\mathrm{Ad} _ {\mathrm{U} _ {s} (\mathrm{G})} a) (x _ {1} x _ {2} \dots x _ {s}). \end{array}
\]

命题 46. 设 \(h \in \mathbf{G}\)、\(x \in \mathrm{T}_h(\mathbf{G})\) 且 \(a \in \mathrm{L}(\mathbf{G})\)。设 \(\phi\) 是映射 \((g, g') \mapsto gg'g^{-1}\)，从 \(\mathbf{G} \times \mathbf{G}\) 到 \(\mathbf{G}\)。像 \(y\) 是 \((a, x) \in \mathrm{T}_e(\mathbf{G}) \times \mathrm{T}_h(\mathbf{G})\) 在 \(\mathrm{T}_{(e,h)}(\phi)\) 下的像，且 \(y = x + h((\mathrm{Ad}h^{-1})a - a)\)。

有

\[
\begin{array}{r l} & y = (\mathrm{T} _ {(e, h)} \phi) (a \otimes \varepsilon_ {h} + \varepsilon_ {e} \otimes x) \\ & \quad = a \top \varepsilon_ {h} + \varepsilon_ {e} \top x \\ & \quad = a * \varepsilon_ {h} - \varepsilon_ {h} * a + x \\ & \quad = h ((\mathrm{Ad} h ^ {- 1}) a) - h a + x. \end{array}
\]

命题 47. 设 G 是李群，H 和 E 是 G 的李子群，并假设 \(hEh^{-1} = E\) 对所有 \(h \in H\) 成立。则 \(\mathcal{T}^{(\infty)}(H) \top \mathcal{T}^{(\infty)}(E) \subset \mathcal{T}^{(\infty)}(E)\)。特别地，\(\operatorname{Ad}(H)(L(E)) \subset L(E)\) 且 \([L(H), L(E)] \subset L(E)\)。

若 \(t \in \mathcal{T}^{(\infty)}(\mathrm{H})\) 且 \(t' \in \mathcal{T}^{(\infty)}(\mathrm{E})\)，则 \(t \otimes t' \in \mathcal{T}^{(\infty)}(\mathrm{H} \times \mathrm{E})\)，并且 \(\mathrm{H} \times \mathrm{E}\) 在映射 \((g, g') \mapsto gg'g^{-1}\) 下的像包含于 \(\mathrm{E}\)。

命题 48. 设 G 是李群，H 和 E 是 G 的李子群。假设作为李群，G 是 H 与 E 的半直积。设 \(\rho\) 是李群 G 的线性表示 \(g \mapsto (\mathrm{Ad} g) \mid \mathrm{L}(E)\)，作用在 \(\mathrm{L}(E)\) 上（参见命题47），并设 \(\sigma\) 是 \(\rho\) 在 H 上的限制。则：

\begin{enumerate}
\def\labelenumi{(\roman{enumi})}
\item
  \(\mathbf{L}(\mathbf{G})\) 是 \(\mathbf{L}(\mathbf{H})\) 与 \(\mathbf{L}(\mathbf{E})\) 的拓扑直和；
\item
  L(H) 是 L(G) 的子代数，而 L(E) 是 L(G) 的理想；
\item
  \(\mathbf{L}(\sigma)\) 是 \(\mathbf{L}(\mathbf{H})\) 在 \(\mathbf{L}(\mathbf{E})\) 的导子李代数上的线性表示；
\item
  L(G) 是由 L(σ) 定义的 L(H) 与 L(E) 的半直积（第 I 章 § 1，第8号）。
\item
  显然，且 (ii) 由命题47推出。\(\mathbf{L}(\sigma) = \mathbf{L}(\rho) \mid \mathbf{L}(\mathbf{H})\)。现在根据命题40（第11号）和命题44（第12号），\(\mathbf{L}(\rho)(t)\) 对所有 \(t \in \mathbf{L}(\mathbf{G})\) 都是 \(\mathrm{ad}_{L(\mathbf{G})} t\) 在 \(\mathbf{L}(\mathbf{E})\) 上的限制。这证明了 (iii)。利用 (i) 和 (ii)，也可证明 (iv)。
\end{enumerate}

推论。设 G 是李群。赋予 \(\mathrm{T}_{e}(\mathrm{G})\) 其唯一的交换李代数结构。设 \(\tau\) 是 \(\mathrm{L}(\mathrm{G})\) 的一个伴随表示。则 \(\mathrm{T}(\mathrm{G})\) 的李代数是 \(\mathrm{L}(\mathrm{G})\) 与 \(\mathrm{T}_{e}(\mathrm{G})\) 的半直积，由 \(\tau\) 定义。换言之，对于 \(x, x'\) 属于 \(\mathrm{L}(\mathrm{G})\) 且 \(y, y'\) 属于 \(\mathrm{T}_{e}(\mathrm{G})\)，

\[
[ (x, y), (x ^ {\prime}, y ^ {\prime}) ] = ([ x, x ^ {\prime} ], [ x, y ^ {\prime} ] + [ y, x ^ {\prime} ])
\]

（左侧的括号在 \(\mathbf{L}(\mathbf{T}(\mathbf{G}))\) 中计算，右侧的括号在 \(\mathbf{L}(\mathbf{G})\) 中计算）。

这由命题48及 § 2 第2号的命题6推出。

命题 49. 设 A 是完备可赋范含幺结合代数。将 A 与 L(A*) 识别。则若 \(g \in \mathbf{A}^*\) 且 \(y \in \mathbf{A}\)，(Ad g) \(y = gyg^{-1}\)。

回忆 \(\operatorname{Ad} g = \mathrm{T}_{1}(\operatorname{Int} g)\)。设 \(u_{g}\) 是从 A 到 A 的映射 \(x \mapsto gxg^{-1}\)。\(\mathrm{A}^{*}\) 到 A 的恒等图将 \(\operatorname{Int} g\) 变为 \(u_{g} \mid \mathrm{A}^{*}\)。这个映射在 \(\mathrm{A}^{*}\) 每一点处的切映射都等于 \(u_{g}\)，故得命题。

推论。对任意 \(g \in \mathbf{A}^*\)，设 \(i(g)\) 是由 \(y \mapsto \mathrm{gyg}^{-1}\) 定义的 \(\mathbf{A}\) 的自同构，从而 \(i\) 是 \(\mathbf{A}^*\) 在 \(\mathbf{A}\) 上的解析线性表示。对任意 \(z \in \mathrm{L}(\mathbf{A}^*) = \mathbf{A}\)，\(\mathrm{L}(i)z\) 是由 \(y \mapsto zy - yz\) 定义的 \(\mathbf{A}\) 的内导子。

这由命题49和命题44推出。

\subsection{张量与不变形式}

设 G 是李群。我们把 G 看作通过左（分别为右）平移作用于自身。设 \(\lambda\) 是对同构为 \(C^{\omega}\) 类的向量函子。则 \(\lambda(\mathrm{TG})\) 是解析的左（分别为右）向量 G-丛（§1，第8号，命题16的推论）。映射 \((g, u) \mapsto gu\)（分别为 ug）从 \(G \times \lambda(\mathrm{L}(G))\) 到 \(\lambda(\mathrm{TG})\) 上是向量 G-丛的同构 \(\phi\)（分别为 \(\psi\)）（§7，第8号，命题17的推论2）。\(\lambda(\mathrm{TG})\) 的每个 G-不变截面都是解析的，并由其在 e 处的值确定（§1，第8号，命题17的推论1）。这样的截面称为左（分别为右）不变。设 \(\sigma\) 是 \(\lambda(\mathrm{TG})\) 的左不变截面；变换 \(\sigma'\) 是 \(\sigma\) 在右平移 \(\delta(g)\) 下的变换，定义为 \(\sigma'(\delta(g)h) = \lambda(\mathrm{T}_{h}(\delta(g)))\sigma(h)\)，对所有 \(h \in G\)；它也是左不变的；它也由 \(\sigma\) 经 \(\gamma(g) \circ \delta(g) = \operatorname{Int}(g)\) 导出，因而

\[
\sigma^ {\prime} (e) = \lambda (\mathrm{Ad} g). \sigma (e).\tag{26}
\]

类似地，设 \(\tau\) 是 \(\lambda(\mathrm{TG})\) 的右不变截面；变换 \(\tau'\) 是 \(\tau\) 在左平移 \(\gamma(g)\) 下的变换，也具有右变性，并且

\[
\tau^ {\prime} (e) = \lambda (\mathrm{Ad} g). \tau (e).\tag{27}
\]

现在考虑 G × G 按下式在左侧作用于 G：

\[
\left(\left(g, g ^ {\prime}\right), g ^ {\prime \prime}\right) \mapsto g g ^ {\prime \prime} g ^ {\prime - 1}.
\]

于是 G 是 \(G \times G\) 的左李齐性空间（§ 1，第6号，例）。因此 \(\lambda(\mathrm{TG})\) 是解析的左向量 \((\mathrm{G} \times \mathrm{G})\)-丛。若 \(\lambda(\mathrm{TG})\) 的一个截面在 \(\mathrm{G} \times \mathrm{G}\) 对 \(\lambda(\mathrm{TG})\) 的作用下不变，或者说在左、右平移下都不变，则称其为双不变截面。设 \(\lambda(\mathrm{L}(\mathrm{G}))_{0}\) 是 \(\lambda(\mathrm{L}(\mathrm{G}))\) 中在 \(\lambda(\mathrm{Ad}(\mathrm{G}))\) 下不变的元素集合。对任意 \(u \in \lambda(\mathrm{L}(\mathrm{G}))_{0}\)，设 \(\sigma_{u}\) 是从 G 到 \(\lambda(\mathrm{TG})\) 的映射，由 \(\sigma_{u}(g) = gu = ug\) 定义。则 \(u \mapsto \sigma_{u}\) 是从 \(\lambda(\mathrm{L}(\mathrm{G}))_{0}\) 到 \(\lambda(\mathrm{TG})\) 的双不变截面集合的双射（§ 1，第8号，命题17的推论1）。

命题 50. 设 G 是李群（当 K 的特征 \textgreater0 时假定其有限维）。设 E 是 \(T_{e}(G)\) 上 k 阶连续交替多重线性形式的向量空间。对任意 \(u \in E\)，设 \(\omega^{u}\) 是 G 上满足如下条件的 k 阶微分形式：\((\omega^{u})_{g}\) 是 \(T_{e}(G)\) 上由平移 \(h \mapsto gh\)（分别为 \(h \mapsto hg\)）从 u 导出的多重线性形式。则 \(\omega^{u}\) 在 G 上解析且左（分别为右）不变。映射 \(u \mapsto \omega^{u}\) 是 E 到 G 上 k 阶左（分别为右）不变微分形式向量空间的同构。

这是上文所述内容的特例。

设 \(\mathbf{F}\) 是完备可赋范空间。若把 \(\mathbf{G}\) 上取值于 \(\mathbf{K}\) 的微分形式替换为 \(\mathbf{G}\) 上取值于 \(\mathbf{F}\) 的微分形式，命题50仍成立。对于每个连续线性映射 \(u\)，它从 \(\mathrm{T}_e(\mathrm{G})\) 映到 \(\mathbf{F}\)，都存在一个 1 阶微分形式 \(\omega^u\)，定义在 \(\mathbf{G}\) 上且取值于 \(\mathbf{F}\)，使得 \((\omega^{u})_{g} = u \circ \mathrm{T}_{g}(\gamma(g)^{-1})\)。特别地，取 \(\mathrm{F} = \mathrm{T}_e(\mathrm{G})\) 及 \(u = \mathrm{Id}_{\mathrm{T_e(G)}}\)。于是得到 G 上的微分形式 \(\omega\)，满足 \(\omega_{g} = \mathrm{T}_{g}(\gamma (g^{-1}))\)；这个微分形式左不变且解析；称为 G 的左典范微分形式。满足 \(\omega_{g}(t) = g^{-1}t\)，对所有 \(t\in \mathrm{T}_g(\mathrm{G})\)。

若 F 再次为任意完备可赋范空间且 \(u \in \mathcal{L}(\mathrm{T}_{e}(\mathrm{G}), \mathrm{F})\)，则 \(\omega^{u} = u \circ \omega\)。特别地（取 F = K），映射 \(v \mapsto v \circ \omega\) 是从 \(\mathrm{T}_{e}(\mathrm{G})\) 的对偶到取值于 K 且在 G 下左不变的 1 阶微分形式向量空间的线性双射。

类似地，G 上的微分形式 \(\omega'\) 满足 \(\omega_g' = \mathrm{T}_g(\delta(g))\)，称为 G 的右典范微分形式。它具有与 \(\omega\) 类似的性质，留给读者陈述。G 到 G 的映射 \(g \mapsto g^{-1}\) 将 \(\omega\) 变换为 \(\omega'\)。

\subsection{莫雷尔-嘉当公式}

设 X 是 \(C^{r}\) 类流形（当 K 的特征 \textgreater0 时为有限维），L 是完备可赋范李代数。设 \(\alpha\) 是 X 上取值于 L 的 \(C^{r-1}\) 类 1 阶微分形式。设 \(x \in X\)。映射

\[
(u _ {1}, u _ {2}) \mapsto [ \alpha_ {x} (u _ {1}), \alpha_ {x} (u _ {2}) ]
\]

从 \(\mathrm{T}_x(\mathbf{X})\times \mathrm{T}_x(\mathbf{X})\) 到 L 的映射是 \(\mathrm{T}_x(\mathbf{X})\) 上取值于 L 的连续交替双线性形式。我们将其记为 \([\alpha ]_x^2\)，于是 \([\alpha ]^2\) 是 X 上取值于 L 的 2 阶微分形式。将 X 中 \(x\) 的一个开邻域识别为 Banach 空间的一个开子集，立即看出 \([\alpha ]^2\) 属于 \(\mathrm{C}^{r - 1}\) 类。若 \(\mathbf{X}'\) 是 \(\mathrm{C}^r\) 类流形且 \(f\colon \mathbf{X}'\to \mathbf{X}\) 是态射，则

\[
[ f ^ {*} (\alpha) ] ^ {2} = f ^ {*} ([ \alpha ] ^ {2}).\tag{28}
\]

设 \(\alpha, \beta\) 是 X 上取值于 L 的两个 \(C^{r-1}\) 类 1 阶微分形式。外积 \(\alpha \wedge \beta\)（《可微与解析流形》，R，7.8.2）是由 \(\alpha\) 与 \(\beta\) 构成的 X 上取值于 L 的 \(C^{r-1}\) 类 2 阶微分形式；有

\[
(\alpha \wedge \beta) _ {x} (u _ {1}, u _ {2}) = [ \alpha_ {x} (u _ {1}), \beta_ {x} (u _ {2}) ] - [ \alpha_ {x} (u _ {2}), \beta_ {x} (u _ {1}) ]\tag{29}
\]

对 \(u_{1}, u_{2}\) 属于 \(\mathrm{T}_x(\mathbf{X})\) 的情形，立即得到

\begin{enumerate}
\def\labelenumi{(\arabic{enumi})}
\setcounter{enumi}{29}
\tightlist
\item
\end{enumerate}

\[
[ \alpha + \beta ] ^ {2} = [ \alpha ] ^ {2} + [ \beta ] ^ {2} + \alpha \wedge \beta\tag{31}
\]

\[
\alpha \wedge \alpha = 2 [ \alpha ] ^ {2}.
\]

命题 51. 设 G 是李群（当 K 的特征 \textgreater0 时为有限维），\(a_{1},\ldots,a_{p}\) 是 L(G) 中的元素，F 是完备可赋范空间，\(\alpha\) 是 G 上取值于 F 的 \(p-1\) 阶微分形式。若 \(\alpha\) 左不变，则

\[
(d \alpha) _ {e} (a _ {1}, \dots , a _ {p}) = \sum_ {i <   j} (- 1) ^ {i + j} \alpha_ {e} ([ a _ {i}, a _ {j} ], a _ {1}, \dots , a _ {i - 1}, a _ {i + 1}, \dots , a _ {j - 1},
\]

\[
\left. a _ {j + 1}, \dots , a _ {p}\right).
\]

若 \(\alpha\) 右不变，则

\[
\begin{array}{l} (d \alpha) _ {e} (a _ {1}, \dots , a _ {p}) \\ = - \sum_ {i <   j} (- 1) ^ {i + j} \alpha_ {e} ([ a _ {i}, a _ {j} ], a _ {1}, \dots , a _ {i - 1}, a _ {i + 1}, \dots , a _ {j - 1}, a _ {j + 1}, \dots , a _ {p}). \end{array}
\]

假设 \(\alpha\) 左不变。根据《可微与解析流形》，R，8.5.7，有

\[
\begin{array}{l} (d \alpha) (L _ {a _ {1}}, \ldots , L _ {a _ {p}}) = \sum_ {i} (- 1) ^ {i - 1} L _ {a _ {i}} \alpha (L _ {a _ {1}}, \ldots , L _ {a _ {i - 1}}, L _ {a _ {i + 1}}, \ldots , L _ {a _ {p}}) \\ + \sum_ {i <   j} (- 1) ^ {i + j} \alpha ([ L _ {a _ {i}}, L _ {a _ {j}} ], L _ {a _ {1}}, \ldots , L _ {a _ {i - 1}}, L _ {a _ {i + 1}}, \ldots , L _ {a _ {j - 1}}, L _ {a _ {j + 1}}, \ldots , L _ {a _ {p}}). \end{array}
\]

但是 G 上的函数 \(\alpha (\mathrm{L}_{a_1},\dots ,\mathrm{L}_{a_{i - 1}},\mathrm{L}_{a_{i + 1}},\dots ,\mathrm{L}_{a_p})\) 左不变，因而为常数。所以

\[
\mathrm{L} _ {a _ {i}} \alpha (\mathrm{L} _ {a _ {1}}, \dots , \mathrm{L} _ {a _ {i - 1}}, \mathrm{L} _ {a _ {i + 1}}, \dots , \mathrm{L} _ {a _ {p}}) = 0.
\]

此外，\([L_{a_{i}}, L_{a_{j}}] = L_{\{a_{i}, a_{j}\}}\)（命题23），从而得到命题51的第一个公式。第二个公式可类似地建立，此时使用关系 \([R_{a_{i}}, R_{a_{j}}] = -R_{\{a_{i}, a_{j}\}}\)。

推论 1. 设 G 是李群（当 K 的特征 \textgreater0 时为有限维），且 ω、ω' 是 G 的左、右典范微分形式。则

\[
d \omega + [ \omega ] ^ {2} = 0, \quad d \omega^ {\prime} - [ \omega^ {\prime} ] ^ {2} = 0.
\]

根据命题51，

\[
\begin{array}{r l} (d \omega) _ {e} (a _ {1}, a _ {2}) & = - \omega_ {e} ([ a _ {1}, a _ {2} ]) = - [ a _ {1}, a _ {2} ] = - [ \omega_ {e} (a _ {1}), \omega_ {e} (a _ {2}) ] \\ & = - [ \omega ] _ {e} ^ {2} (a _ {1}, a _ {2}) \end{array}
\]

从而得到第一个公式。第二个公式可类似地建立。

推论 2. 假设 G 是有限维的。设 \((e_1, \ldots, e_n)\) 是 L(G) 的一组基，\((e_1^*, \ldots, e_n^*)\) 是对偶基，\((c_{ijk})\) 是相对于基 \((e_1, \ldots, e_n)\) 的 L(G) 的结构常数，并设 \(\omega_i\)（分别为 \(\omega_i^{\prime})\)）是 G 上取值于 K 的左（分别为右）不变微分形式，满足 \((\omega_i)_e = e_i^*\)（分别为 \((\omega_i^{\prime})_e = e_i^*)\)。则

\[
d \omega_ {k} + \sum_ {i <   j} c _ {i j k} \omega_ {i} \wedge \omega_ {j} = 0 (k = 1, 2, \dots , n)
\]

\[
d \omega_ {k} ^ {\prime} - \sum_ {i <   j} c _ {i j k} \omega_ {i} ^ {\prime} \wedge \omega_ {j} ^ {\prime} = 0 (k = 1, 2, \dots , n).
\]

若 \(r < s\)，

\[
\begin{array}{r l} \left(d \omega_ {k}\right) _ {e} (e _ {r}, e _ {s}) & = - \left(\omega_ {k}\right) _ {e} ([ e _ {r}, e _ {s} ]) \\ & = - \sum_ {i} c _ {r s i} (\omega_ {k}) _ {e} (e _ {i}) \\ & = - c _ {r s k} \\ & = - \sum_ {i <   j} c _ {i j k} (\omega_ {i} \wedge \omega_ {j}) _ {e} (e _ {r}, e _ {s}). \end{array}
\]

对于 \(\omega_{k}^{\prime}\)，论证类似。

\subsection{不变微分形式的构造}

引理 2. 设 \(\mathbf{G}\) 是李群，\(\mathrm{U}\) 是 \(e\) 在 \(\mathbf{G}\) 中的一个对称开邻域，\(\mathrm{E}\) 是完备可赋范空间，且 \(\phi: \mathrm{U}^2 \to \mathrm{E}\) 是解析映射。对任意 \(g \in \mathrm{U}\)，设 \(\omega_g\) 是某映射在点 \(g\) 处的微分，该映射为 \(h \mapsto \phi(g^{-1}h)\)。则 \(\omega\) 是 \(\mathrm{U}\) 上的左不变微分形式，它是 \(\mathbf{G}\) 上在 \(e\) 处的值为 \(d_e\phi\) 的左不变微分形式的限制。

显然 \(\omega_{e} = d_{e}\phi\)。对任意 \(g \in \mathrm{U}\) 及 \(t \in \mathrm{T}_{e}(\mathrm{G})\)，

\[
\begin{array}{r l} \langle \omega_ {g}, \mathrm{T} _ {e} (\gamma (g)) t \rangle & = \langle d _ {g} (\phi \circ \gamma (g) ^ {- 1}), \mathrm{T} _ {e} (\gamma (g)) t \rangle \\ & = \langle d _ {e} \phi \circ \mathrm{T} _ {g} (\gamma (g) ^ {- 1}), \mathrm{T} _ {e} (\gamma (g)) t \rangle = \langle d _ {e} \phi , t \rangle \end{array}
\]

因此 \(\omega_{g}\) 由 \(d_{e}\phi\) 经 \(\mathrm{T}_{e}(\gamma(g))\) 导出。

命题 52. 设 \(n\) 是整数 \(>0\)，\(G\) 是 \(n\) 维李群，\(U\) 是 \(e\) 在 \(G\) 中的对称开邻域，且 \(\psi: U^2 \to K^n\) 是 \(G\) 的一个图，满足 \(\psi(e) = 0\)。若 \((x_1, \ldots, x_n)\) 是 \(x \in \psi(U)\) 的坐标，\((y_1, \ldots, y_n)\) 是 \(y \in \psi(U)\) 的坐标，我们记

\[
m _ {1} (x _ {1}, \dots , x _ {n}, y _ {1}, \dots , y _ {n}), \dots , m _ {n} (x _ {1}, \dots , x _ {n}, y _ {1}, \dots , y _ {n})
\]

\(\psi (\psi^{-1}(x)^{-1}\psi^{-1}(y))\) 的坐标。于是，若对 \(1\leqslant k\leqslant n\) 写成

\[
\begin{array}{r l} \varpi_ {k} (x _ {1}, \ldots , x _ {n}) = D _ {n + 1} m _ {k} (x _ {1}, \ldots , x _ {n}, x _ {1}, \ldots , x _ {n}) d x _ {1} + \dots \\ & \quad + \mathrm{D} _ {2 n} m _ {k} (x _ {1}, \ldots , x _ {n}, x _ {1}, \ldots , x _ {n}) d x _ {n}, \end{array}\tag{32}
\]

则微分形式 \(\varpi_k\) 定义在 \(\psi(\mathbf{U})\) 上，由 \(\psi\) 从 \(\mathbf{G}\) 上的左不变微分形式导出，并且满足 \(\varpi_k(0,\ldots,0) = dx_k\)。

应用引理2，取 E = K，并令 \(\phi(g)\) 为 \(\psi(g)\) 的指标为 k 的坐标。得到微分形式 \(\omega_{k}\)；设 \(\varpi_{k}\) 是它在 \(\psi\) 下的变换。\(\varpi_{k}\) 在 \((x_{1},\ldots,x_{n})\) 处的值，是在 \((x_{1},\ldots,x_{n})\) 处对函数 \(y\mapsto m_{k}(x_{1},\ldots,x_{n},y_{1},\ldots,y_{n})\) 求得的微分；因此该值由公式 (32) 给出。于是只需使用引理2的结论。

命题 53. 设 G 是李群，A 是完备可赋范代数，且 \(\phi\) 是从 G 到 A* 的李群态射。对任意 \(g \in G\)，设 \(\omega_g = \phi(g)^{-1}.d_g\phi\)。则 \(\omega\) 是 G 上在 \(e\) 处取值为 \(d_e\phi\) 的左不变微分形式。

应用引理2，取 \(\mathbf{E} = \mathbf{A}\) 且 \(\mathbf{U} = \mathbf{G}\)。在 \(g\) 处，映射 \(h \mapsto \phi(g^{-1}h) = \phi(g)^{-1}\phi(h)\) 的微分是 \(\phi(g)^{-1}.d_g\phi\)。

\subsection{李群上的 Haar 测度}

设 G 是 n 维李群。于是 \(\bigwedge^{n}(\mathrm{T}_{e}(\mathrm{G}))\) 是 1 维的。因此（第13号），G 上 n 阶左不变微分形式的向量空间 S 是 1 维的。设 \((\omega_{1},\ldots,\omega_{n})\) 是 G 上 1 阶左不变微分形式空间的一组基；则 \(\omega_{1}\wedge\omega_{2}\wedge\cdots\wedge\omega_{n}\) 是 S 的一组基。

命题 54. 设 G 是 n 维李群，\(\omega\) 是 G 上的 n 阶左不变微分形式，\(\phi\) 是 G 的自同态。则

\[
\phi^ {*} (\omega) = (\det L (\phi)) \omega .
\]

记 \(\mathrm{L}(\phi) = u, \omega_e = f\) 且 \(\phi^*(\omega)_e = g\)。对任意 \(x_1, \ldots, x_n\)，它们属于 \(\mathrm{L}(\mathrm{G})\)，

\[
g (x _ {1}, \dots , x _ {n}) = f (u x _ {1}, \dots , u x _ {n}) = (\det u) f (x _ {1}, \dots , x _ {n})
\]

因此 \(\phi^{*}(\omega)_{e} = \det L(\phi). \omega_{e}\)。另一方面，若 \(g \in G\)，

\[
\phi \circ \gamma (g) = \gamma (\phi (g)) \circ \phi
\]

因此 \(\gamma(g)^{*}\phi^{*}(\omega)=\phi^{*}(\omega)\)。所以 \(\phi^{*}(\omega)\) 左不变，命题得证。

推论。对任意 \(g \in G\)，

\[
\delta (g) ^ {*} \omega = (\det \operatorname{Ad} g) \omega .
\]

\[
\delta (g) ^ {*} \omega = \delta (g) ^ {*} \gamma (g) ^ {*} \omega = (\operatorname{Int} g) ^ {*} \omega \text {   and   } L (\operatorname{Int} g) = \operatorname{Ad} g.
\]

设 G 是局部紧群，\(\phi\) 是 G 的自同态。假设存在 e 的开邻域 V、\(V'\)，使得 \(\phi(V) = V'\)，且 \(\phi \mid V\) 是从 G 到 G 的局部同构。设 \(\mu\) 是 G 上的左 Haar 测度。根据《积分学》第 VII 章 §1 命题9的推论，存在唯一的数 a \textgreater{} 0，使得 \(\phi(\mu \mid V) = a^{-1}\mu \mid V'\)。显然，a 与 V、\(V'\) 及 \(\mu\) 的选择无关。称其为 \(\phi\) 的模，记为 mod \(_{G}\) \(\phi\)，或简记为 mod \(\phi\)。当 \(\phi\) 是 G 的自同构时，便得到《积分学》第 VII 章 §1 的定义4。

命题 55. 假设 K 是局部紧的。设 \(\mu\) 是 K 的加法群上的 Haar 测度。设 G 是 \(n\) 维李群。

\begin{enumerate}
\def\labelenumi{(\roman{enumi})}
\item
  设 \(\omega\) 是 G 上非零的 \(n\) 阶左不变微分形式。则测度 mod \((\omega)_{\mu}\)（《可微与解析流形》，R，10.1.6）是 G 上的左 Haar 测度。若 K = R 且 G 具有由 \(\omega\) 定义的定向，则由 \(\omega\) 定义的测度（《可微与解析流形》，R，10.4.3）是 G 上的左 Haar 测度。
\item
  设 \(\phi\) 是 G 的 étale 自同态。则 mod \(\phi = \text{mod}\det L(\phi)\)。
\item
  显然。设 \(\mathrm{V}, \mathrm{V}'\) 是 \(e\) 的开邻域，使得 \(\phi(\mathrm{V}) = \mathrm{V}'\)，且 \(\phi \mid \mathrm{V}\) 是从 \(\mathrm{G}\) 到 \(\mathrm{G}\) 的局部同构。则
\end{enumerate}

\[
\begin{array}{r l r l} \phi^ {- 1} (\mathrm{mod} (\omega) _ {\mu}   |   \mathrm{V} ^ {\prime}) & = \mathrm{mod} (\phi^ {*} (\omega)) _ {\mu}   |   \mathrm{V}) & & \text {by transport of structure} \\ & = \mathrm{mod} (\det \mathrm{L} (\phi) \omega   |   \mathrm{V}) _ {\mu} & & \text {(Proposition 54)} \\ & = \mathrm{mod} \det \mathrm{L} (\phi) (\mathrm{mod} (\omega) _ {\mu}   |   \mathrm{V}) \end{array}
\]

因此 \(\phi = \mathrm{mod}\det \mathrm{L}(\phi)\)；这是由 \(\phi\) 的模的定义得到的。

推论。对任意 \(g \in \mathbf{G}\)，\(\Delta_{\mathbf{G}}(g) = (\text{mod det } \operatorname{Ad} g)^{-1}\)。特别地，\(\mathbf{G}\) 为幺模群的充要条件是 mod det \(\operatorname{Ad} g = 1\) 对所有 \(g \in \mathbf{G}\) 都成立。

\[
\begin{array}{r l} \Delta_ {\mathbf {G}} (g) & = (\text {mod} \operatorname{Int} g) ^ {- 1} \quad \text {(Integration, Chapter VII, § 1, formula (33))} \\ & = (\text {mod} \det \operatorname{L} (\operatorname{Int} g)) ^ {- 1} \quad \text {(Proposition 55)} \\ & = (\text {mod} \det \operatorname{Ad} g) ^ {- 1}. \end{array}
\]

注记。在保持命题52的假设和记号的情况下，假设 K 是局部紧的。设 \(\mu\) 是测度

\[
\operatorname{mod} \det (\mathrm{D} _ {n + i} m _ {k} (x _ {1}, \dots , x _ {n}, x _ {1}, \dots , x _ {n})) _ {1 \leqslant i, k \leqslant n} d x _ {1} \dots d x _ {n}
\]

在 \(\psi (\mathbf{U})\) 上。则 \(\psi^{-1}(\mu)\) 是 G 上某个 Haar 测度限制到 U 的结果。

命题 56. 设 G 是 n 维李群，H 是 p 维李子群，X 是李齐性空间 G/H。假设

\[
\det \operatorname{Ad} _ {\mathrm{L} (\mathrm{G})} h = \det \operatorname{Ad} _ {\mathrm{L} (\mathrm{H})} h
\]

对所有 \(h\in \mathbf{H}\) 均成立。则：

\begin{enumerate}
\def\labelenumi{(\roman{enumi})}
\item
  \(n - p\) 阶的 \(\mathbf{X}\) 上在 \(\mathbf{G}\) 下不变的微分形式都是解析的。
\item
  这些形式构成的向量空间是 1 维的。
\item
  若 \(\omega\) 是这样的非零形式且 K 是局部紧的，则 mod \((\omega)_{\mu}\) 是 X 上的非零测度，并且在 G 下不变。
\end{enumerate}

根据§ 1 第8号的例，\(\mathrm{Alt}^{n - p}(\mathrm{TX},\mathrm{K})\) 是解析向量 G-丛。设 \(x_0\) 是 \(e\) 在 \(\mathbf{X}\) 中的典范像；其稳定子为 H。\(\mathrm{Alt}^{n - p}(\mathrm{TX},\mathrm{K})\) 在 \(x_0\) 处的纤维是 \(\bigwedge^{n - p}\mathrm{T}_{x_0}(\mathbf{X})^*\)，而 \(\mathrm{T}_{x_0}(\mathbf{X})\) 与 \(\mathrm{L}(G) / L(\mathrm{H})\) 典范地识别。若 \(h\in \mathrm{H}\)，\(\tau_h\) 是 \(\mathbf{X}\) 的自同构，由 \(h\) 定义；在取商时，它由 G 的自同构 \(g\mapsto hgh^{-1}\) 导出。因此，自同构 \(\mathrm{T}_{x_0}(\tau_h)\) 在取商时由 \(\mathrm{Ad}_{\mathrm{L(G)}}(h)\) 导出。由于

\[
\det \mathrm{Ad} _ {\mathrm{L(G)}} h = (\det \mathrm{Ad} _ {\mathrm{L(H)}} h). (\det \mathrm{T} _ {x _ {0}} (\tau_ {h})),
\]

假设推出 \(\det \mathrm{T}_{x_0}(\tau_h) = 1\)。所以 \(\bigwedge^{n - p}\mathrm{T}_{x_0}(\mathbf{X})^*\) 的每个元素都在 H 下不变。于是 (i) 和 (ii) 由§ 1 第8号命题17的推论1推出，而 (iii) 显然。

由《积分学》第 VII 章 § 2 定理3的推论2，存在一个在 G 下不变的 X 上非零正测度；因为命题56的假设推出 \(\Delta_{\mathbf{G}}|H = \Delta_{\mathbf{H}}\)（命题55的推论）。

命题 57. 设 G 是 n 维李群。~取 \(\bigwedge^{n}\mathrm{T}_{e}(\mathrm{G})^{*}\) 的一组基；借助 \(\bigwedge^{n}\mathrm{T}(\mathrm{G})^{*}\) 的右（分别为左）平凡化，我们可以将这个向量丛与平凡向量丛 \(\mathbf{G} \times \mathbf{K}\) 识别，从而标量微分算子的转置被识别为标量微分算子。

于是，若 \(u \in \mathrm{U}(\mathrm{G})\)，\(\mathrm{L}_u\)（分别为 \(\mathbf{R}_u\)）的转置是 \(\mathrm{L}_u^\vee\)（分别为 \(\mathbf{R}_u^\vee\)）。

我们考虑将 \(\bigwedge^{n}\mathrm{T}(\mathrm{G})^{*}\) 用右不变形式 \(\omega\) 平凡化的情形。

假设对于 U(G) 中的元素 \(u_{1}, u_{2}\) 命题已经得到证明。则

\[
\begin{array}{l l} ^ {t} (\mathrm{L} _ {u _ {1} * u _ {2}}) = ^ {t} (\mathrm{L} _ {u _ {1}} \circ \mathrm{L} _ {u _ {2}}) & \text {(Proposition 23)} \\ = ^ {t} (\mathrm{L} _ {u _ {2}}) \circ^ {t} (\mathrm{L} _ {u _ {1}}) & \text {(Diff. \& Anal. Man., R, 14.3.3)} \\ = \mathrm{L} _ {u _ {2}} ^ {\vee} \circ \mathrm{L} _ {u _ {1}} ^ {\vee} & \text {by hypothesis} \\ = \mathrm{L} _ {u _ {2} * u _ {1}} ^ {\vee \vee} & \text {(Proposition 23)} \\ = \mathrm{L} _ {(u _ {1} * u _ {2}) \vee} & \text {(Proposition 7)} \end{array}
\]

因此命题对 \(u_{1} * u_{2}\) 成立。所以只需在 \(u \in \mathrm{T}_{e}(\mathrm{G})\) 时证明命题。现在，\(\mathbf{L}_u\) 由 G 在 G 上的右作用定义（第6号），因而 \(\theta_{\mathrm{L}_u}\omega = 0\)，因为 \(\omega\) 右不变（《可微与解析流形》，R，8.4.5）；所以，若 \(f\) 是定义在 \(e\) 的一个开邻域上且取值于 K 的解析函数，则 \(\theta_{\mathrm{L}_u}(f\omega) = (\theta_{\mathrm{L}_u}f)\omega\)（《可微与解析流形》，R，8.4.8）。利用上述识别以及《可微与解析流形》R，14.4.1，\(\mathbf{L}_u\) 的转置是 \(-\mathbf{L}_u\)，即 \(\mathbf{L}_u^\vee\)。

推论。设 G 是有限维实李群，\(\mu\)（分别为 \(\nu\)）是 G 上的左（分别为右）Haar 测度，k 是满足 \(\geqslant0\) 的整数，\(u\in\mathrm{U}_{k}(\mathrm{G})\)，f 和 g 是 G 上具有紧支集的 \(C^{k}\) 类实值函数。则

\[
\int_ {G} \left(R _ {u} f\right) g d \mu = \int_ {G} f (R _ {u} ^ {\vee} g) d \mu
\]

\[
\int_ {\mathbf {G}} \left(\mathrm{L} _ {u} f\right) g d \nu = \int_ {\mathbf {G}} f (\mathrm{L} _ {u} ^ {\vee} g) d \nu .
\]

这由命题57及《可微与解析流形》R，14.3.8推出。

\subsection{左微分}

定义 9. 设 G 是李群，M 是 \(\mathbf{C}^r\) 类流形，\(f\) 是从 M 到 G 的 \(\mathbf{C}^r\) 类映射。\(f\) 的左（分别为右）微分是 M 上取值于 L(G) 的 1 阶微分形式，它把每个向量 \(u \in \mathrm{T}_m(\mathrm{M})\) 对应到元素 \(f(m)^{-1}\) . \((\mathrm{T}_m f)(u)\)（分别为 \((\mathrm{T}_m f)(u). f(m)^{-1}\)）。

本章只考虑左微分，记为 \(f^{-1}.df\)，右微分的相应结果留给读者翻译。

若 \(f\) 是恒等映射；\(G, f^{-1}.df\) 是左典范微分形式 \(\omega\)，属于 \(G\)。回到定义8的一般情形，

\[
(f ^ {- 1}. d f) _ {m} = \omega_ {f (m)} \circ \mathrm{T} _ {m} (f)
\]

因此 \(f^{-1}.df = f^{*}(\omega)\)。这说明 \(f^{-1}.df\) 属于 \(\mathbf{C}^{r - 1}\) 类。

例。(1) 若 G 是完备可赋范空间的加法群，且 \(T_{0}(E)\) 与 E 典范地识别，则 \(f^{-1}\) . df 就是《可微与解析流形》R，8.2.2 中定义的微分 df。

\begin{enumerate}
\def\labelenumi{(\arabic{enumi})}
\setcounter{enumi}{1}
\tightlist
\item
  假设 G 是与完备可赋范代数 A 相关的乘法群 A*。于是 f 可看作从 M 到 A 的映射，因而《可微与解析流形》R，8.2.2 意义下的微分 df 有定义，而《可微与解析流形》R，8.3.2 意义下的乘积 \(f^{-1}df\) 也有定义。显然，后一个形式与 f 的左微分相同。
\end{enumerate}

命题 58. 设 G 和 H 是两个李群，M 是 \(\mathbf{C}^r\) 类流形，\(f\) 是从 M 到 G 的 \(\mathbf{C}^r\) 类映射，\(h\) 是从 G 到 H 的态射。则

\[
(h \circ f) ^ {- 1}. d (h \circ f) = \mathrm{L} (h) \circ (f ^ {- 1}. d f) = (h ^ {- 1}. d h) \circ \mathrm{T} (f).
\]

对任意 \(x \in \mathbf{M}\) 及 \(u \in \mathrm{T}_x(\mathrm{M})\)，

\[
(h \circ f) ^ {- 1}. d (h \circ f) (u) = ((h \circ f) (x)) ^ {- 1}. T (h \circ f) (u).
\]

后一表达式一方面等于

\[
\begin{array}{l l} \mathrm{T} (h) (f (x) ^ {- 1}. \mathrm{T} (f) (u)) & (\S 2, \text { Proposition   5 }) \\ = \mathrm{T} _ {e} (h) ((f ^ {- 1}. d f) (u)) \end{array}
\]

另一方面等于

\[
\begin{array}{c} h (f (x)) ^ {- 1} \mathrm{T} (h) (\mathrm{T} (f) u) \\ = (h ^ {- 1}. d h) (\mathrm{T} (f) u). \end{array}
\]

命题 59. 设 G 是李群，M 是 \(\mathbf{C}^r\) 类流形，\(f\) 和 \(g\) 是从 M 到 G 的 \(\mathbf{C}^r\) 类映射，\(p\) 是从 TM 到 M 的典范满射。

\[
(i) (f g) ^ {- 1}. d (f g) = (\mathrm{Ad} \circ g \circ p) ^ {- 1} \circ (f ^ {- 1}. d f) + g ^ {- 1}. d g.
\]

\begin{enumerate}
\def\labelenumi{(\roman{enumi})}
\setcounter{enumi}{1}
\tightlist
\item
  定义 \(h(m) = f(m)^{-1}\)，其中 \(m \in M\) 任意。
\end{enumerate}

\[
h ^ {- 1}. d h = - (\mathrm{Ad} \circ f \circ p) \circ (f ^ {- 1}. d f).
\]

断言 (i) 由 § 2 第2号命题7推出。令 \(g = h\)，由 (i) 得到断言 (ii)。

推论 1. 设 \(s \in G\)，且 \(sg\) 是映射 \(x \mapsto sg(x)\)，其定义域为 \(M\)，陪域为 \(G\)。则 \((sg)^{-1}.d(sg) = g^{-1}.dg\)。

这由命题59 (i) 推出，其中取 \(f\) 为常值映射 \(x \mapsto s\)，其定义域为 \(M\)，陪域为 \(G\)。

推论 2. 若映射 \(f\) 和 \(g\) 的定义域为 \(M\)、陪域为 \(G\)，且它们具有相同的左微分，则 \(fg^{-1}\) 的切映射处处为零。若进一步 \(K\) 的特征为 0，则 \(fg^{-1}\) 局部为常值。

根据命题59，

\[
(f g ^ {- 1}) ^ {- 1}. d (f g ^ {- 1}) = (\mathrm{Ad} \circ g \circ p) \circ (f ^ {- 1}. d f) - (\mathrm{Ad} \circ g \circ p) \circ (g ^ {- 1}. d g).
\]

若 \(f^{-1}.df = g^{-1}.dg\)，则 \((fg^{-1})^{-1}.d(fg^{-1}) = 0\)，即 \(T_x(fg^{-1}) = 0\) 对所有 \(x \in M\) 成立。这证明了第一个断言。第二个断言由《可微与解析流形》R，5.5.3从它推出。

命题 60. 设 G 是李群（当 K 的特征 \textgreater0 时为有限维），M 是 \(\mathbf{C}^r\) 类流形，\(f\) 是从 M 到 G 的 \(\mathbf{C}^r\) 类映射，\(\alpha\) 是 \(f\) 的左微分。则 \(d\alpha + [\alpha]^2 = 0\)。

设 \(\omega\) 是 G 的左典范微分形式。利用第14号命题51的推论1，有

\[
\begin{array}{r l} d \alpha & = d (f ^ {*} (\omega)) = f ^ {*} (d \omega) = f ^ {*} (- [ \omega ] ^ {2}) \\ & = - [ f ^ {*} (\omega) ] ^ {2} = - [ \alpha ] ^ {2}. \end{array}
\]

\subsection{李群芽的李代数}

本号中，\((G, e, \theta, m)\) 表示一个李群芽。本节的大部分结果仍以同样的证明成立。我们复述其中将要用到的结果。

18.1. 设 \(\Omega\) 是 \(m\) 的定义域。设 \((g, g') \in \Omega\)、\(t \in \mathrm{T}_g^{(\alpha)}(\mathrm{G})\)、\(t' \in \mathrm{T}_{g'}^{(\alpha)}(\mathrm{G})\)。如第1号所述，\(t\) 与 \(t'\) 的卷积乘积记为 \(t * t'\)，它是 \(t \otimes t'\) 在 \(m\) 下的像。记 \(\mathrm{U}(\mathrm{G}) = \mathrm{T}_e^{(\alpha)}(\mathrm{G})\)、\(\mathrm{U}_s(\mathrm{G}) = \mathrm{T}_e^{(s)}(\mathrm{G})\)、\(\mathrm{U}^+(\mathrm{G}) = \mathrm{T}_e^{(\alpha)+}(\mathrm{G})\)、\(\mathrm{U}_s^+(\mathrm{G}) = \mathrm{T}_e^{(s)+}(\mathrm{G})\)。对于 \(t\) 和 \(t'\)（它们属于 \(\mathrm{U}(\mathrm{G})\)），\(t * t'\) 有定义且属于 \(\mathrm{U}(\mathrm{G})\)。在卷积乘积下，\(\mathrm{U}(\mathrm{G})\) 是单位元为 \(\varepsilon_e\) 的结合代数，并由 \(\mathrm{U}_s(\mathrm{G})\) 滤化。典范同构 \(i_{\mathrm{G}, e}\) 将 gr \(\mathrm{U}(\mathrm{G})\) 映到 \(\mathrm{TS}(\mathrm{T}_e(\mathrm{G}))\) 上，并且是代数同构。

18.2. 设 G、H 是李群芽，\(\phi: G \to H\) 是一个态射。若 \(t \in \mathrm{U}(\mathrm{G})\)，则 \(\mathrm{U}(\phi)(t)\) 是 t 在 \(\phi_{*}\) 下的像，是 U(H) 的元素，且 \(\mathrm{U}(\phi)\) 是从代数 U(G) 到代数 U(H) 的态射。G 到 G 的映射 \(\theta: x \mapsto x^{-1}\) 定义了 U(G) 到 U(G) 的映射 \(t \mapsto t^{\vee}\)。对于 U(G) 中的 \(t, t'\)，\(t * t'\) 相对于 \(G^{\vee}\) 计算时，等于 \(t' * t\) 相对于 G 计算时的结果，且 \((t * t')^{\vee} = t^{\prime\vee} * t^\vee\)。于是 \(\mathrm{U}(\phi)(t^{\vee}) = (\mathrm{U}(\phi)t)^{\vee}\)。若 \(G_{1}, \ldots, G_{n}\) 是李群芽且 \(G = G_{1} \times \cdots \times G_{n}\)，则从 \(\mathrm{U}(G_{1}) \otimes \cdots \otimes \mathrm{U}(G_{n})\) 到 U(G) 的典范同构是代数同构；

对 U(G) 中的 \(t_{1},\ldots,t_{n}\)，\((t_{1}\otimes\cdots\otimes t_{n})^{\vee}=t_{1}^{\vee}\otimes\cdots\otimes t_{n}^{\vee}\)。设 H 是 G 的李子群芽，i: H→G 是典范嵌入。则 U(i) 是从代数 U(H) 到 U(G) 的单射同态，并且

\[
\mathrm{U} (i) \left(t ^ {\vee}\right) = \left(\mathrm{U} (i) (t)\right) ^ {\vee}
\]

对所有 \(t \in \mathrm{U}(\mathrm{H})\) 均成立。在由 G 的流形结构定义的卷积乘积和余乘法下，\(\mathrm{U}(\mathrm{G})\) 是双代数，\(\mathrm{U}(\phi)\) 是双代数态射。

18.3. 设 G 是李群芽，X 是 \(C^{r}\) 类流形，\(\psi\) 是 G 在 X 上的 \(C^{r}\) 类左作用律片段。设 \(\Omega\) 是 \(\psi\) 的定义域。若 \(t \in T_{g}^{(s)}(G)\)、\(u \in T_{x}^{(s^{\prime\prime})}(X)\)、\((g, x) \in \Omega\) 且 \(s + s^{\prime} \leqslant r\)，则令 \(t * u\) 表示 \(t \otimes u\) 在 \(\psi_{*}\) 下的像。设 \(t \in T_{g}^{(s)}(G)\)、\(t^{\prime} \in T_{g'}^{(s^{\prime})}(G)\)、\(u \in T_{x}^{(s^{\prime})}(X)\)；若 \(s + s^{\prime} + s^{\prime\prime} \leqslant r\) 且 \(gg^{\prime}\)、\((gg^{\prime})x\)、\(g^{\prime}x\)、\(g(g^{\prime}x)\) 有定义，则

\[
(t * t ^ {\prime}) * u = t * (t ^ {\prime} * u).
\]

设 \(x_0 \in \mathbf{X}\)，\(\rho(x_0)\) 是映射 \(g \mapsto gx_0\)，它在 \(e\) 的某个开邻域上有定义。若 \(t \in \mathrm{U}_r(\mathrm{G})\)，则 \(\rho(x_0)_*t = t * \varepsilon_{x_0}\)。这里以及本号其余部分，右作用律片段的相应结果留给读者翻译。

18.4. 保持 18.3 的记号，设 \(t \in \mathrm{U}_s(\mathbf{G})\) 且 \(s \leqslant r\)。设 \(f\) 是 \(\mathbf{C}^r\) 类函数，定义在 \(\mathbf{X}\) 上且取值于 Hausdorff 多赋范空间。令 \(t * f\) 表示由下式定义的 \(\mathbf{X}\) 上的函数：

\[
\begin{array}{r l} (t * f) (x) & = \langle t, g \mapsto f (\psi (\theta (g), x)) \rangle \\ & = \langle t ^ {\vee}, f \circ \rho (x) \rangle = \langle \rho (x) _ {*} (t ^ {\vee}), f \rangle = \langle t ^ {\vee} * \varepsilon_ {x}, f \rangle . \end{array}
\]

若 \(t \in \mathrm{U}_s(\mathrm{G})\)、\(t' \in \mathrm{U}_{s'}(\mathrm{G})\) 且 \(s + s' \leqslant r\)，则 \(\langle t', t * f \rangle = \langle t^\vee * t', f \rangle\) 且 \((t * t') * f = t * (t' * f)\)。设 \(t \in \mathrm{U}_s(\mathrm{G})\)，\(f\) 和 \(f'\) 是 \(C^r\) 类函数，定义在 X 上且分别取值于 Hausdorff 多赋范空间 \(F, F'\)，\((u, u') \mapsto u. u'\) 是从 \(F \times F'\) 到 Hausdorff 多赋范空间的连续双线性映射；令 \(\underline{n}\)。

\(\sum_{i=1}^{n} t_i \otimes t_i'\) 是 \(t\) 在余乘法下的像；若 \(s \leqslant r\)，则

\[
t * (f f ^ {\prime}) = \sum_ {i = 1} ^ {n} (t _ {i} * f) (t _ {i} ^ {\prime} * f ^ {\prime}).
\]

18.5. 保持 18.3 的记号，设 \(t \in \mathrm{U}_s(\mathrm{G})\) 且 \(s \leqslant r\)。映射 \(x \to t * \varepsilon_x\) 称为由 \(t\) 和作用律片段定义的点分布场，有时记为 \(\mathrm{D}_t^{\psi}\) 或 \(\mathrm{D}_t\)。若 \(f: \mathrm{X} \to \mathrm{F}\) 是 \(C^r\) 类函数，则函数 \(t^{\vee} * f\) 在 \(\mathrm{X}\) 上也记为 \(\mathrm{D}_t f\)；它属于 \(C^{r - s}\) 类（当 \(s < \infty\) 时）。若 \(t \in \mathrm{U}_s(\mathrm{G})\)、\(t' \in \mathrm{U}_{s'}(\mathrm{G})\) 且 \(s + s' \leqslant r\)，则 \(\mathrm{D}_{t,t'} f = \mathrm{D}_{t'}(\mathrm{D}_t f)\)。若 \(G\) 和 \(X\) 是有限维的，则 \(D_t\) 是 \(X\) 上阶数 \(\leqslant s\) 且属于 \(C^{r - s}\) 类的微分算子（若 \(s < \infty\)）。函数 \(D_t f\) 就是 \(f\) 在这个微分算子下的变换。

18.6. 设 G 是李群芽且 \(t \in \mathrm{U}(\mathrm{G})\)。\(L_{t}\) 表示 G 上的点分布场 \(g \mapsto \varepsilon_{g} * t\)，而 \(R_{t}\) 表示点分布场 \(g \mapsto t * \varepsilon_{g}\)。若 \(f \in \mathcal{C}^{\omega}(\mathrm{G}, \mathrm{F})\)，则 \(\mathrm{L}_{t} f \in \mathcal{C}^{\omega}(\mathrm{G}, \mathrm{F})\) 且 \(\mathrm{R}_{t} f \in \mathcal{C}^{\omega}(\mathrm{G}, \mathrm{F})\)。对于 \(t, t'\)，它们属于 \(\mathrm{U}(\mathrm{G})\)，\(L_{t* t'} = L_t \circ L_{t'}, R_{t* t'} = R_{t'} \circ R_t, L_t \circ R_{t'} = R_{t'} \circ L_t, \theta(L_t) = R_t^\vee\)。

18.7. 由于 \(T_{e}(G)\) 是 U(G) 的本原元素集合，

\[
[ \mathrm{T} _ {e} (\mathrm{G}), \mathrm{T} _ {e} (\mathrm{G}) ] \subset \mathrm{T} _ {e} (\mathrm{G}).
\]

可赋范空间 \(\mathrm{T}_{e}(\mathrm{G})\) 与括号一起构成可赋范李代数，称为 G 的可赋范李代数（或 G 的李代数），记为 \(\mathrm{L}(\mathrm{G})\)。设 \(\mathrm{E}(\mathrm{G})\) 是 \(\mathrm{L}(\mathrm{G})\) 的包络代数。\(\mathrm{L}(\mathrm{G})\) 到 \(\mathrm{U}(\mathrm{G})\) 的典范嵌入定义了同态 \(\eta\)，它从代数 \(\mathrm{E}(\mathrm{G})\) 到代数 \(\mathrm{U}(\mathrm{G})\)；若 K 的特征为 0，则 \(\eta\) 是双代数同构，借此将 \(\mathrm{U}(\mathrm{G})\) 与 \(\mathrm{E}(\mathrm{G})\) 识别。采用 18.3 的记号，对所有 \(a \in \mathrm{L}(\mathrm{G})\)，设 \(D_{a}\) 是由 a 在 X 上定义的点分布场。映射 \((a, x) \mapsto D_{a}(x)\) 是 \(C^{r-1}\) 类态射，从平凡向量丛 \(\mathrm{L}(\mathrm{G}) \times \mathrm{X}\) 到向量丛 \(\mathrm{T}(\mathrm{X})\)。设 I 是包含 0 的 K 的一个开子集，\(\gamma: I \to G\) 是 \(C^{r}\) 类映射，满足 \(\gamma(0) = e\)。令 \(a = T_{0}(\gamma) 1 \in L(\mathrm{G})\)。若 f: X → F 是 \(C^{r}\) 类函数，则

\[
\left(\mathrm{D} _ {a} f\right) (x) = \lim _ {k \in \mathrm{K} ^ {*}, k \rightarrow 0} k ^ {- 1} \left(f (\gamma (k) x) - f (x)\right).
\]

若 \(r \geqslant 2\)，则映射 \(a \mapsto \mathrm{D}_a\) 是 \(\mathrm{C}^{r-1}\) 类的 \(\mathrm{L}(\mathrm{G})\) 在 \(\mathrm{X}\) 上的左无穷小作用律。

18.8. 设 G 和 H 是李群芽，\(\phi\) 是从 G 到 H 的态射。\(\mathrm{U}(\phi)\) 在 \(\mathrm{L}(\mathrm{G})\) 上的限制，也就是 \(\mathrm{T}_{e}(\phi)\)，是从 \(\mathrm{L}(\mathrm{G})\) 到 \(\mathrm{L}(\mathrm{H})\) 的连续态射，记为 \(\mathrm{L}(\phi)\)。若 \(\psi\) 是从 H 到某个李群芽的态射，则 \(\mathrm{L}(\psi \circ \phi) = \mathrm{L}(\psi) \circ \mathrm{L}(\phi)\)。\(\phi\) 是浸入的充要条件是，\(\mathrm{L}(\phi)\) 是从 \(\mathrm{L}(\mathrm{G})\) 到 \(\mathrm{L}(\mathrm{H})\) 中某个具有拓扑补空间的李子代数的同构。特别地，若 G 是 H 的李子群芽且 \(\phi\) 是典范嵌入，则将 \(\mathrm{L}(\mathrm{G})\) 与 \(\mathrm{L}(\mathrm{H})\) 的一个李子代数借助 \(\mathrm{L}(\phi)\) 识别。若 \((\mathrm{G}_{i})_{i\in I}\) 是李群芽的有限族且 G 是它们的乘积，则 \(\mathrm{L}(\mathrm{G})\) 与 \(\prod_{i\in I}\mathrm{L}(\mathrm{G}_{i})\) 典范地识别。

18.9. 设 G 是李群芽（当 K 的特征 \textgreater0 时为有限维）。设 F 是完备可赋范空间。设 \(\alpha\) 是 G 上取值于 F 的 k 阶微分形式。称 \(\alpha\) 在 G 上左不变，当 \(\alpha_{g}\) 由 \(\alpha_{e}\) 经映射 \(h \mapsto gh\) 从 e 的某个邻域导出到 g 的某个邻域时。若 \(\alpha\) 左不变，则 \(\alpha\) 解析。映射 \(\alpha \mapsto \alpha_{e}\) 是从 G 上取值于 F 的 k 阶左不变微分形式集合到从 \(T_{e}(G)\) 到 F 的连续交替 k-线性映射集合的双射。若 \(\alpha_{e} = \operatorname{Id}_{\mathrm{T}_e(\mathrm{G})}\)，则称 \(\alpha\) 为 G 的左典范微分形式。右不变微分形式及右典范微分形式的定义类似。若 \(\omega\) 是 G 的左典范微分形式，则 \(d\omega + [\omega]^{2} = 0\)。设 M 是 \(C^{r}\) 类流形，f 是从 M 到 G 的 \(C^{r}\) 类映射。f 的左微分记为 \(f^{-1}.df\)，它是 M 上取值于 L(G) 的 1 阶微分形式，把每个向量 \(u \in T_{m}(M)\) 对应到元素 \(f(m)^{-1}\) . \((\mathrm{T}_{m}f)(u)\)。则 \(f^{-1}.df = f^{*}(\omega)\) 且 \(d\alpha + [\alpha]^{2} = 0\)。若从 M 到 G 的两个映射 f 和 g 具有相同的左微分且 K 的特征为 0，则 \(fg^{-1}\) 局部为常值。

\section{从李代数到李群的过渡}

回忆一下，在本章结尾之前，均假设 K 的特征为 0。

\subsection{从李代数态射到李群态射的过渡}

引理 1. 设 \(\mathbf{G}\) 是李群芽，\(\mathfrak{h}\) 是 \(\mathbf{L}(\mathbf{G})\) 中具有拓扑补空间的李子代数。\(\mathfrak{gh}\)（分别为 \(\mathfrak{hg}\)）对 \(g\in\mathbf{G}\) 的并集是 \(\mathbf{T}(\mathbf{G})\) 的可积向量子丛。

考虑 \(\mathrm{T}(\mathrm{G})\) 的左平凡化（§ 2，第3号），立即看出，\(g\mathfrak{h}\) 对于 \(g\in \mathrm{G}\) 是向量子丛 \(\mathrm{E}\)（它是 \(\mathrm{T}(\mathrm{G})\) 的子丛）的纤维。设 \(g\in \mathrm{G}\)。\((\mathrm{L}_a)_g\)（其中 \(a\in \mathfrak{h}\)）的集合等于 \(g\mathfrak{h}\)。现在，若 \(a\) 和 \(b\) 属于 \(\mathfrak{h}\)，则 \([\mathrm{L}_a,\mathrm{L}_b] = \mathrm{L}_{[a,b)}\) 且 \([a,b]\in \mathfrak{h}\)。因此 \(\mathrm{E}\) 可积（《可微与解析流形》，R，9.3.3 (iv)）。对于 \(\mathfrak{h}g\)，论证类似。

\(\mathfrak{gh}\)（分别为 \(\mathfrak{hg}\)）的并集的积分叶状结构（《可微与解析流形》，R，9.3.2）称为与 \(\mathfrak{h}\) 相关的 G 的左（分别为右）叶状结构。

定理 1. 设 \(\mathbf{G}\) 和 \(\mathbf{H}\) 是李群芽，\(f\) 是从 \(\mathbf{L}(\mathbf{G})\) 到 \(\mathbf{L}(\mathbf{H})\) 的连续态射。

\begin{enumerate}
\def\labelenumi{(\roman{enumi})}
\item
  存在一个开李子群芽 \(\mathbf{G}'\)，它是 \(\mathbf{G}\) 的开李子群芽，以及一个态射 \(\phi\)，其定义域为 \(\mathbf{G}'\)、陪域为 \(\mathbf{H}\)，使得 \(f = \mathbf{L}(\phi)\)。
\item
  设 \(\mathrm{G}_1, \mathrm{G}_2\) 是 G 的开李子群芽，\(\phi_i\) 是从 \(\mathrm{G}_i\) 到 H 的态射，满足 \(f = \mathrm{L}(\phi_i)\)，其中 \(i = 1, 2\)。则 \(\phi_1\) 与 \(\phi_2\) 在 \(e\) 的某个邻域上重合。设 \(p_1: \mathrm{G} \times \mathrm{H} \to \mathrm{G}\)、\(p_2: \mathrm{G} \times \mathrm{H} \to \mathrm{H}\) 是典范投影。对所有 \((g,h) \in \mathrm{G} \times \mathrm{H}\)，设 \(f_{g,h}\) 是映射 \(ga \mapsto hf(a)\)，从 \(\mathrm{T_g(G)} = g\mathrm{L(G)}\) 到 \(\mathrm{T_h(H)} = h\mathrm{L(H)}\)。考虑 \(\mathrm{T(G)}\) 和 \(\mathrm{T(H)}\) 的左平凡化，立即看出 \(f_{g,h}\) 定义了从 \(p_1^*\mathrm{T(G)}\) 到 \(p_{2}^{*}\mathbf{T}(\mathrm{H})\) 的态射。设 \(a\) 是 \(f\) 的图；它是 \(\mathbf{L}(\mathrm{G}) \times \mathbf{L}(\mathrm{H})\) 的闭李子代数，并以 \(\{0\} \times \mathbf{L}(\mathrm{H})\) 为拓扑补空间。对所有 \((g,h) \in \mathrm{G} \times \mathrm{H}\)，\(f_{g,h}\) 的图是 \((g,h)\) . \(a\)。这些图的并集是 \(\mathrm{T}(\mathrm{G} \times \mathrm{H})\) 的可积向量子丛（引理1）。于是存在（《可微与解析流形》，R，9.3.7）一个开邻域 \(\mathrm{U}\)，包含 \(e_{\mathrm{G}}\) 且位于 \(\mathrm{G}\) 中，以及一个解析映射 \(\phi\)，其定义域为 \(\mathrm{U}\)、陪域为 \(\mathrm{H}\)，使得 \(\phi(e_{\mathrm{G}}) = e_{\mathrm{H}}\)，并且 \(\mathrm{T}_{g}(\phi) = f_{g,\phi(g)}\) 对所有 \(g \in \mathrm{U}\)。特别地，\(\mathrm{T}_{e_{\mathrm{G}}}(\phi) = f\)。
\end{enumerate}

设 \(\mathbf{V}\) 是包含 \(e_{\mathrm{G}}\) 且位于 \(\mathbf{G}\) 中的开邻域，使得对 \((s,t)\in \mathbf{V}\times \mathbf{V}\)，乘积 \(st\) 和 \(\phi (s)\phi (t)\) 有定义且 \(st\in \mathbf{U}\)。考虑映射 \(\alpha_{1},\alpha_{2}\)，它们定义在 \(\mathbf{V}\times \mathbf{V}\) 上、取值于 \(\mathbf{H}\)，并由下式给出：

\[
\alpha_ {1} (s, t) = \phi (t s), \quad \alpha_ {2} (s, t) = \phi (t) \phi (s).
\]

于是 \(\alpha_{1}(t,e) = \phi (t) = \alpha_{2}(t,e)\)。另一方面，固定 V 中的 \(t\)，设 \(\beta_{i}\) 是从 V 到 H 的映射 \(s\mapsto \alpha_i(s,t)\)。则对所有 \(s\in \mathrm{V}\) 及 \(a\in \mathrm{L}(\mathrm{G})\)，

\[
\begin{array}{r l} \mathrm{T} _ {s} (\beta_ {1}) (s a) & = \mathrm{T} _ {t s} (\phi) (t s a) = f _ {t s, \phi (t s)} (t s a) \\ & = \phi (t s) f (a) = f _ {s, \beta_ {1} (s)} (s a) \\ \mathrm{T} _ {s} (\beta_ {2}) (s a) & = \phi (t) \mathrm{T} _ {s} (\phi) (s a) = \phi (t) f _ {s, \phi (s)} (s a) \\ & = \phi (t) \phi (s) f (a) = f _ {s, \beta_ {2} (s)} (s a). \end{array}
\]

因此，根据《可微与解析流形》R，9.3.7，\(\alpha_{1}\) 与 \(\alpha_{2}\) 在 \((e_{\mathrm{G}}, e_{\mathrm{G}})\) 的某个邻域上重合。因此，将 \(\phi\) 限制在 \(e_{\mathrm{G}}\) 的充分小的对称开邻域上，得到李群芽的态射，从而得到 (i)。

设 \(\mathrm{G}_1, \mathrm{G}_2, \phi_1, \phi_2\) 如 (ii) 中，并证明 \(\phi_1, \phi_2\) 在 \(e_{\mathrm{G}}\) 的某个邻域上重合。存在 \(e_{\mathrm{G}}\) 的一个开邻域 W，使得 \(\phi_1(t s) = \phi_1(t)\phi_1(s)\)、\(\phi_2(t s) = \phi_2(t)\phi_2(s)\) 对其中所有 \(s\)、\(t\) 成立。于是，若 \(s \in W\) 且 \(a \in L(G)\)，

\[
\mathrm{T} _ {s} \left(\phi_ {i}\right) (s a) = \phi_ {i} (s) \mathrm{T} _ {e} \left(\phi_ {i}\right) (a) = \phi_ {i} (s) f (a) = f _ {s, \phi _ {i} (s)} (s a)
\]

对 \(i = 1,2\) 均成立。由于 \(\phi_1(e_{\mathrm{G}}) = e_{\mathrm{H}} = \phi_2(e_{\mathrm{G}})\)，由《可微与解析流形》R，9.3.7，\(\phi_1\) 与 \(\phi_2\) 在 \(e_{\mathrm{G}}\) 的某个邻域上重合。

推论 1. 设 G 和 H 是两个李群芽。若 L(G) 与 L(H) 同构，则 G 与 H 局部同构。

这由定理1及 § 1 第10号命题21推出。

推论 2. 设 \(\mathbf{G}\) 是李群芽。若 \(\mathrm{L}(\mathbf{G})\) 是交换的，则 \(\mathbf{G}\) 与加法李群 \(\mathrm{L}(\mathbf{G})\) 局部同构。

加法群 \(\mathbf{L}(\mathbf{G})\) 的李代数与 \(\mathbf{L}(\mathbf{G})\) 同构。因此只需应用推论1。

推论 3. 设 G 是李群。若 L(G) 是交换的，则 G 含有一个交换开子群。

存在 G 的一个交换开李子群芽 U（推论2）。设 V 是 e 的一个邻域，满足 \(V^{2} \subset U\)。则对 V 中的所有 x、y，都有 xy = yx。因此，由 V 生成的 G 的子群是交换的；它显然是开的。

\subsection{从李代数到李群的过渡}

我们将 Hausdorff 级数记为 H(X, Y)（第 II 章 § 6，第4号，定义1）。

引理 2. 设 \(\mathbf{L}\) 是定义在 \(\mathbf{R}\) 或 \(\mathbf{C}\) 上的完备赋范李代数。设 \(\mathbf{G}\) 是 \(x \in \mathbf{L}\) 的集合，其中 \(\| x \| < \frac{1}{3} \log \frac{3}{2}\)。设 \(\theta\) 是映射 \(x \mapsto -x\)，从 \(\mathbf{G}\) 到 \(\mathbf{G}\)。设 \(\mathbf{H}\) 是限制在 \(\mathbf{G} \times \mathbf{G}\) 上的 \(\mathbf{L}\) 的 Hausdorff 函数（第 II 章 § 7，第2号）。

\begin{enumerate}
\def\labelenumi{(\roman{enumi})}
\item
  (G, 0, θ, H) 是李群芽。
\item
  设 \(\phi\) 是从 G 到 L 的恒等映射。\(\phi\) 在 0 处的微分是从可赋范李代数 L(G) 到 L 的同构。
\end{enumerate}
(i) 由第 II 章 § 7 第2号推出。

由于 \(\phi\) 是 G 上的图，在 0 处的微分 \(\psi\) 是 \(\phi\) 给出的可赋范空间同构。另一方面，映射 H 的积分级数展开 \(\mathrm{H} = \sum_{i,j\geqslant 0}\mathrm{H}_{ij}\) 满足 \(\mathrm{H}_{11}(x,y) = \frac{1}{2} [x,y]\)。根据§ 3 命题24，对任意 \(a,b\)，它们属于 \(\mathrm{L}(\mathrm{G})\)，

\[
\psi ([ a, b ]) = \mathrm{H} _ {1 1} (\psi (a), \psi (b)) - \mathrm{H} _ {1 1} (\psi (b), \psi (a)) = [ \psi (a), \psi (b) ]
\]

这证明了 (ii)。

称 G 为由 L 定义的李群芽。

假设 K 是超度量的。设 p 是 K 的剩余域的特征。若 \(p \neq 0\)，令 \(\lambda = |p|^{1/(p-1)}\)；若 p = 0，令 \(\lambda = 1\)。

引理 3. 设 L 是定义在 K 上的完备赋范李代数。设 G 是 \(x \in \mathbf{L}\) 的集合，满足 \(\| x \| < \lambda\)。设 H: G × G → G 是 L 的 Hausdorff 函数（第 II 章 § 8，第3号）。

\begin{enumerate}
\def\labelenumi{(\roman{enumi})}
\item
  在复合律 H 下，G 是一个李群，其中 0 是单位元，对所有 \(x \in G\)，-x 是 x 的逆元。
\item
  设 \(\phi\) 是从 \(\mathbf{G}\) 到 \(\mathbf{L}\) 的恒等映射。\(\phi\) 在 0 处的微分是从可赋范李代数 \(\mathrm{L}(\mathrm{G})\) 到 \(\mathbf{L}\) 的同构。
\item
  对任意 \(\mu \in \mathbf{R}_+^*\)，设 \(\mathrm{G}_{\mu}\) 是 \(x \in \mathbf{L}\) 的集合，满足 \(\| x \| < \mu\)。则这些 \(\mathrm{G}_{\mu}\) 在 \(\mu < \lambda\) 时构成 0 的开闭邻域基本系，并且是 \(\mathbf{G}\) 的子群。
\end{enumerate}

断言 (i) 和 (iii) 由第 II 章 § 8 第3号命题3推出，而 (ii) 可像引理2中那样证明。

称 G 为由 L 定义的李群。

定理 2. 设 \(\mathbf{L}\) 是完备可赋范李代数。存在一个李群芽 \(G\)，使得 \(\mathrm{L}(\mathbf{G})\) 与 \(\mathbf{L}\) 同构。任意两个这样的李群芽局部同构。

第一个断言由引理2和引理3推出。第二个断言由第1号定理1的推论1推出。

推论 1. 设 \(\mathbf{G}\) 是李群。存在一个 \(e\) 的邻域，其中不含异于 \(\{e\}\) 的有限子群。若 \(\mathbf{K} = \mathbf{R}\) 或 \(\mathbf{C}\)，则存在一个 \(e\) 的开邻域，其中不含异于 \(\{e\}\) 的子群。

记 \(\mathrm{L}(G) = \mathrm{L}\)。在 \(\mathbf{L}\) 上取一个定义 \(\mathbf{L}\) 拓扑的范数，使得 \(\| [x,y]\| \leqslant \| x\| \| y\|\) 对所有 \(x,y\) 属于 \(\mathbf{L}\) 均成立。

假设 \(\mathbf{K} = \mathbf{R}\) 或 \(\mathbf{C}\)。设 \(G'\) 是由 L 定义的李群芽。存在一个以 0 为中心的开球 \(U'\)，位于 \(G'\) 中，以及一个同构 \(\phi\)，其定义域是李群芽 \(U'\)，陪域为开邻域 \(U\)，其中包含 \(e\)，且位于 G 中。令 \(V' = \frac{1}{2} U'\)、\(V = \phi(V')\)，设 H 是包含于 V 的 G 的子群，且 \(h \in H\)。令 \(x = \phi^{-1}(h) \in V'\)。若 \(x \neq 0\)，则存在整数 \(n > 0\)，使得 \(x, 2x, \ldots, nx\) 属于 \(V'\)，\((n + 1)x \in U'\)，而 \((n + 1)x \notin V'\)。于是 \(h, h^2, \ldots, h^n\) 属于 V，

\[
h ^ {n + 1} \in \mathrm{U}, \quad h ^ {n + 1} \notin \mathrm{V},
\]

这不可能。因此 \(\mathbf{H} = \{e\}\)。

假设 K 是超度量的。只需在 G 是与 L 相关的李群时证明此推论。若 \(g \in G\)，则 g 在 G 中计算的幂，是 Zg 在 L 中计算的元素。当 \(g \neq e\) 时这些元素彼此不同。因此 G 不含异于 \(\{e\}\) 的有限子群。

推论 2. 设 \(k\) 是 \(\mathbf{K}\) 的非离散闭子域，\(\mathbf{G}\) 是定义在 \(k\) 上的李群，且 \(\mathbf{L} = \mathbf{L}(\mathbf{G})\)。假设 \(\mathbf{L}\) 具有可赋范李 \(\mathbf{K}\)-代数结构 \(\mathbf{L}'\)，它与可赋范李 \(k\)-代数结构相容，并且在 \(\mathbf{G}\) 的伴随表示下不变。则 \(\mathbf{G}\) 上存在唯一一个李 \(\mathbf{K}\)-群结构，它与李 \(k\)-群结构相容，且其李代数为 \(\mathbf{L}'\)。

存在定义在 K 上的李群芽 \(G_{1}\)，使得 \(\mathrm{L}(G_{1}) = \mathrm{L}'\)（定理2）。根据第1号定理1的推论1，将 G 和 \(G_{1}\) 看作李 k-群芽时，它们局部同构。因此，G 中存在 e 的一个开邻域 \(G'\)，以及 \(G'\) 上的李 K-群芽结构，其李代数为 L，并且与定义在 k 上的李群芽结构相容。设 V 是 G 中 e 的对称开邻域，满足 \(V^{2} \subset G'\)。设 \(g \in G\)。则 \(\phi = Int g\) 是从 \(G'\) 的某个充分小的开李子群芽到 \(G'\) 的一个开李子群芽的 k-同构；并且 \(T_{e}(\phi)\) 是 K-线性的，因而当 x 充分接近 e 时 \(T_{x}(\phi)\) 是 K-线性的；所以 Int g 限制在 V 中 e 的某个充分小的开邻域上是 K-解析的（《可微与解析流形》，R，5.14.6）。根据§1第9号命题18，G 上存在一种解析 K-流形结构，使 G 成为一个李

K-群，且 V 是 G 的开子-K-流形。通过平移可见，G 的底层 k-流形结构就是给定的结构。李 K-群 G 的李代数与开李 K-子群芽 V 的李代数相同，因而为 L'。最后，推论中所述的唯一性由 § 3 第8号命题32推出。

定理 3. 设 G 是李群芽，\(\mathfrak{h}\) 是 \(\mathrm{L}(\mathrm{G})\) 中具有拓扑补空间的李子代数。存在 G 的李子群芽 H，使得 \(\mathrm{L}(\mathrm{H}) = \mathfrak{h}\)。若 \(\mathrm{H}_1\) 和 \(\mathrm{H}_2\) 是 G 的李子群芽，且

\[
\mathrm{L} (\mathrm{H} _ {1}) = \mathrm{L} (\mathrm{H} _ {2}) = \mathfrak {h},
\]

则 \(\mathrm{H}_1\cap \mathrm{H}_2\) 在 \(\mathrm{H}_1\) 和 \(\mathrm{H}_2\) 中都是开的。

存在一个李群芽 H'，其李代数与 \(\mathfrak{h}\) 同构（定理2）。必要时缩小 H'，可设存在从 H' 到 G 的态射 \(\phi\)，使得 \(\mathrm{L}(\phi)\) 是从 \(\mathrm{L}(\mathrm{H}')\) 到 \(\mathfrak{h}\) 的同构（第1号定理1）。由于 \(\mathfrak{h}\) 具有拓扑补空间，\(\phi\) 在 e 处是浸入。因此进一步缩小 H'，可设 \(\phi\) 是从流形 H' 到 G 的一个子流形的同构。这证明了 H 的存在性。第二个断言由下列命题推出：

命题 1. 设 G 是李群芽，H 和 H' 是两个李子群芽。要使 L(H) ⊃ L(H')，充要条件是 H ∩ H' 在 H' 中开。若 H ∩ H' 在 H' 中开，则 L(H') = L(H ∩ H') ⊂ L(H)。假设 L(H) ⊃ L(H')。设 i、i' 是 H、H' 到 G 的典范嵌入。必要时缩小 H'，可设存在从 H' 到 H 的态射 ψ，使得 L(ψ) 是 L(H') 到 L(H) 的典范嵌入（第1号定理1）。于是 L(i ∘ ψ) = L(i')，因此存在 H' 中 eH 的一个邻域 V，使得 i ∘ ψ 与 i' 在 V 上重合（定理1）。所以 V ⊂ H，进而 V ⊂ H ∩ H'，且 H ∩ H' 在 H' 中开（§ 1，第10号）。

命题 2. 设 G 是定义在 K 上的李群，k 是 K 的非离散闭子域，H 是李 k-群 G 的李子群。假设 L(H) 是 L(G) 的一个具有拓扑补空间的向量子-K-空间。则 H 是李 K-群 G 的李子群。

存在李 K-群 G 的一个李子群芽 \(\mathrm{H}'\)，使得 \(\mathrm{L}(\mathrm{H}') = \mathrm{L}(\mathrm{H})\)（定理3）。将 G、H、H' 看作李 \(k\)-群芽；定理3于是证明 \(\mathrm{H} \cap \mathrm{H}'\) 在 H 和 H' 中开。因此存在 G 中 e 的一个开邻域 U，使得 \(e\) 的 \(\mathrm{U} \cap \mathrm{H}\) 是 G 上的 K-子流形。所以 H 是李 K-群 G 的李子群（\(\S 1\)，第3号，命题6）。

\subsection{指数映射}

定理 4. 设 G 是李群芽，L 是其李代数，V 是 L 中 0 的一个开邻域，\(\phi\) 是从 V 到 G 的解析映射，满足 \(\phi(0) = 0\) 及 \(\mathrm{T}_0(\phi) = \mathrm{Id}_L\)。以下条件等价：

\begin{enumerate}
\def\labelenumi{(\roman{enumi})}
\item
  对所有 \(b \in \mathbf{L}\)，有 \(\phi((\lambda + \lambda')b) = \phi(\lambda b)\phi(\lambda'b)\)，其中 \(|\lambda|\) 和 \(|\lambda'|\) 充分小。
\item
  对所有 \(b \in \mathbf{L}\) 及每个整数 \(n > 0\)，\(\phi_{*}(b^{n})\) 是 \(n\) 次齐次元素，位于

\(\mathrm{U}(\mathrm{G})\) 中（将 \(\mathrm{T}_{0}^{(\infty)}(\mathrm{L})\) 与 \(\mathrm{TS}(\mathrm{L})\) 识别，并把 \(b^{n}\) 视为在 \(\mathrm{TS}(\mathrm{L})\) 中计算）。

\item
  TS(L) 到 U(G) 的映射 \(\phi_{*}\) 与 TS(L) 和 U(G) 的分次相容。
\item
  TS(L) 到 U(G) 的映射 \(\phi_*\) 是 TS(L) 到 L 的包络代数的典范映射。
\item
  存在一个定义 \(\mathbf{L}\) 的拓扑的范数，并且使 \(\mathbf{L}\) 满足
\end{enumerate}

\[
\| [ x, y ] \| \leqslant \| x \| \| y \|
\]

对所有 \(x, y\)，其中元素属于 \(\mathbf{L}\)，并且存在一个开子群芽 \(\mathbf{W} \subset \mathbf{V}\)，它是由 \(\mathbf{L}\) 定义的李群芽（第2号）的开子群芽，使得 \(\phi | \mathbf{W}\) 是从 \(\mathbf{W}\) 到 \(\mathbf{G}\) 的一个开李子群芽的同构。

(v) \(\Rightarrow\) (i)：显然，因为在 W 中有 \((\lambda b).(\lambda^{\prime}b) = (\lambda +\lambda^{\prime})b\)，其中 \(|\lambda |\) 和 \(|\lambda '|\) 充分小。
(i) \(\Rightarrow\) (ii)：假设条件 (i) 成立。设 \(b \in L\)。设 \(\psi\) 是 \(\phi\) 限制在 \(V \cap Kb\) 上的映射。根据假设，存在 0 的一个对称邻域 \(T\)，它位于加法李群 \(Kb\) 中，使得 \(\psi|T\) 是从李群芽 \(T\) 到 \(G\) 的态射。因此

\[
\phi_ {*} (b ^ {n}) = (\psi | \mathrm{T}) _ {*} (b ^ {n}) = ((\psi | \mathrm{T}) _ {*} (b)) ^ {n} = (\phi_ {*} (b)) ^ {n},
\]

所以 \(\phi_{*}(b^{n})\) 是 \(n\) 次齐次元素，属于 U(G)。

(ii) \(\Rightarrow\) (iii)：这是因为 \(\mathsf{TS}^n (\mathbf{L})\) 是 \(\mathsf{TS}(\mathbf{L})\) 的向量子空间，由 \(n\) 次幂（取自 \(\mathbf{L}\) 中元素）生成（《代数》第 IV 章 § 5 命题5）。
(iii) \(\Rightarrow\) (iv)：TS(L) 到 L 的包络代数的典范映射，是将 1 映到 1 且延拓 \(\mathrm{Id}_{\mathrm{L}}\) 的分次余代数的唯一态射（第 II 章 § 1，第5号，注记3）。现在，根据假设，\(\phi_{*}\) 是余代数态射且 \(\phi_{*}|_{\mathrm{L}} = \mathrm{Id}_{\mathrm{L}}\)。若条件 (iii) 成立，容易看出条件 (iv) 也成立。
(iv) \(\Rightarrow\) (v)：假设条件 (iv) 成立。在 L 上取一个定义 L 的拓扑的范数，使 \(\| [x,y]\| \leqslant \| x\| \| y\|\) 对所有 \(x,y\) 成立。设 H 是由赋范李代数 L 定义的李群芽。根据定理1，存在 H 的一个开子群芽 \(S\subset V\) 以及从 S 到 G 的一个开子群芽的同构 \(\phi^{\prime}\)。由于我们已经知道 (v) \(\Rightarrow\) (iv)，TS(L) 到 U(G) 的映射 \(\phi_{*}\) 是 TS(L) 到 L 的包络代数的典范映射。因此有 \(\phi_{*}(t) = \phi_{*}^{\prime}(t)\)，对所有 \(t\in T_0^{(\infty)}(L)\)。由于 \(\phi\) 和 \(\phi^{\prime}\) 都解析，\(\phi\) 与 \(\phi^{\prime}\) 在 0 的某个邻域上重合。

定义 1. 设 G 是李群芽，L 是其李代数。G 的指数映射是定义在 L 中 0 的一个开邻域上、取值于 G 且满足定理4条件的任意解析映射 \(\phi\)。

定理4立即推出：对每个李群芽 G，都存在 G 的指数映射，并且 G 的两个指数映射在 0 的某个邻域上重合。

例。(1) 设 G 是完备可赋范空间 E 的加法群。L(G) 到 E 的典范同构满足定理4的条件 (i)，因而是 G 的指数映射。

\begin{enumerate}
\def\labelenumi{(\arabic{enumi})}
\setcounter{enumi}{1}
\tightlist
\item
  设 A 是完备赋范含幺结合代数。设 \(\mathbf{A}^*\) 是由 A 的可逆元素组成的李群。我们将 \(\mathrm{L}(\mathrm{A}^*)\) 与 A 识别（§3，第9号，命题33的推论）。若 \(\mathbf{K} = \mathbf{R}\) 或 \(\mathbf{C}\)，我们知道第 II 章 §7 第3号中定义的从 A 到 \(\mathbf{A}^*\) 的映射 exp 满足定理4的条件 (i)，因而是指数映射。现在设 K 是超度量的。设 \(p\) 是 K 的剩余域的特征。若 \(p \neq 0\)，令 \(\lambda = |p|^{1/(p-1)}\)；若 \(p = 0\)，令 \(\lambda = 1\)。设 U 是 \(x \in \mathbf{A}\) 且满足 \(\|x\| < \lambda\) 的集合。我们知道（第 II 章 §8，第4号），从 U 到 \(\mathbf{A}^*\) 的映射 exp 满足定理4的条件 (i)，因而是指数映射。注意 U 是 A 的加法子群。
\end{enumerate}

这个例子说明了定义1所采用术语的由来。

设 \(G\) 是李群芽，\(\phi\) 是 \(G\) 的指数映射。则 \(\phi\) 在 0 处是 étale 的，因而存在一个开邻域 \(U\)，它位于 \(L(G)\) 中的 0 附近，使得 \(\phi(U)\) 在 \(G\) 中开，且 \(\phi|U\) 是从解析流形 \(U\) 到解析流形 \(\phi(U)\) 的同构。

G 上的典范图（第一类）是解析流形 G 上的一个图 \(\psi\)，其逆映射是指数映射。若进一步 G 是有限维的，并选定 L(G) 的一组基，则由 \(\psi\) 和该基在 \(\psi\) 的定义域上定义的坐标系称为典范坐标系（第一类）。

命题 3. 设 \(\mathbf{G}\) 是李群芽，\(\mathbf{L}\) 是其李代数，\(\phi\) 是 \(\mathbf{G}\) 的指数映射。设 \(\mathbf{L}_1, \ldots, \mathbf{L}_n\) 是 \(\mathbf{L}\) 的向量子空间，使得 \(\mathbf{L}\) 是 \(\mathbf{L}_1, \ldots, \mathbf{L}_n\) 的拓扑直和。映射

\[
(b _ {1}, b _ {2}, \dots , b _ {n}) \mapsto \theta (b _ {1}, b _ {2}, \dots , b _ {n}) = \phi (b _ {1}) \phi (b _ {2}) \dots \phi (b _ {n}),
\]

定义在 \(L_{1} \times L_{2} \times \cdots \times L_{n}\) 的一个开子集上，是解析的。在 \((0, 0, \ldots, 0)\) 处，\(\theta\) 的切映射是从 \(L_{1} \times \cdots \times L_{n}\) 到 L 的典范映射。

设 \(k_{i}\) 是从 \(L_{i}\) 到 \(L_{1} \times L_{2} \times \cdots \times L_{n}\) 的典范嵌入。则对所有 \(b \in L_{i}\)，\((\mathrm{T}_{(0,\ldots,0)}\theta)(\mathrm{T}_{0}\phi_{i})(b) = (\mathrm{T}_{0}\phi)(b) = b\)，从而 \((\mathrm{T}_{(0,\ldots,0)}\theta)|L_{i}\) 是从 \(L_{i}\) 到 L 的典范嵌入。

特别地，\(\theta\) 在 \((0, 0, \ldots, 0)\) 处是 étale 的。它限制在 \((0,0,\ldots,0)\) 的充分小开邻域 U 上时，在 G 中的像 \(\theta(\mathrm{U})\) 是开的，且这个限制是一个同构。逆映射 \(\eta\) 把 \(\theta(\mathrm{U})\) 映到 U，称为 G 的第二类典范图，它与给定的 L 的直和分解相关。若进一步 G 是有限维的，且每个 \(L_{i}\) 都由非零向量 \(e_{i}\) 生成，则在 \(\theta(\mathrm{U})\) 上由 \(\eta\) 和 \(e_{i}\) 定义的坐标系称为第二类典范坐标系。

命题 4. 设 G 是李群芽，\(\phi\) 是 G 的单射指数映射。对 L(G) 中的所有 \(x, y\)，


\[
x + y = \lim _ {\lambda \in \mathbb {K} ^ {*}, \lambda \rightarrow 0} \lambda^ {- 1} \phi^ {- 1} (\phi (\lambda x) \phi (\lambda y))\tag{1}
\]

\[
[ x, y ] = \lim _ {\lambda \in \mathbb {K} ^ {*}, \lambda \rightarrow 0} \lambda^ {- 2} \phi^ {- 1} (\phi (\lambda x) \phi (\lambda y) \phi (- \lambda x) \phi (- \lambda y))\tag{2}
\]

（注意，\(\phi^{-1}(\phi(\lambda x)\phi(\lambda y))\) 和 \(\phi^{-1}(\phi(\lambda x)\phi(\lambda y)\phi(-\lambda x)\phi(-\lambda y))\) 在 \(|\lambda|\) 充分小时有定义。）

给 \(\mathbf{L} = \mathbf{L}(\mathbf{G})\) 赋予一个定义 \(\mathbf{L}\) 拓扑的范数，使 \(\| [x,y] \| \leqslant \| x \| \| y \|\) 对所有 \(x, y\) 属于 \(\mathbf{L}\) 成立。利用定理2和定理4，可以假设 \(\mathbf{G}\) 是由 \(\mathbf{L}\) 定义的李群芽，且 \(\phi = \mathrm{Id}_{\mathbf{G}}\)。记 \((x,y) \mapsto x.y\) 为群 \(\mathbf{G}\) 中的乘积。于是待证公式可写为


\[
x + y = \lim _ {\lambda \in \mathbf {K} ^ {*}, \lambda \rightarrow 0} \lambda^ {- 1} ((\lambda x). (\lambda y))\tag{3}
\]

\[
[ x, y ] = \lim _ {\lambda \in \mathbb {K} ^ {*}, \lambda \rightarrow 0} \lambda^ {- 2} ((\lambda x). (\lambda y). (- \lambda x). (- \lambda y)).\tag{4}
\]

存在 K 中 0 的一个开邻域 V，使函数

\[
\lambda \mapsto f (\lambda) = (\lambda x). (\lambda y)
\]

在 V 上有定义且解析。根据第 II 章 § 6 第4号注记2，\(f\) 在原点处的积分级数展开为

\[
\lambda (x + y) + \frac {1}{2} \lambda^ {2} [ x, y ] + \dots
\]

这证明了 (3)。另一方面，对于 G 中的 \(u, v\)，当 \(\| u \|\)、\(\| v \|\) 充分小时，\(u.v\) 是 \((u, v)\) 的解析函数，而这个函数在原点处的积分级数展开中的 1 次和 2 次项为 \(u + v + \frac{1}{4}[u, v]\)。根据《可微与解析流形》R，3.2.7 和 4.2.3，函数 \(f(\lambda)f(-\lambda)\) 在原点处的积分级数展开中的 1 次和 2 次项，就是

\[
f (\lambda) + f (- \lambda) + \frac {1}{2} [ f (\lambda), f (- \lambda) ]
\]

或者也可写为

\[
\begin{array}{r l} \lambda (x + y) + \frac {1}{2} \lambda^ {2} [ x, y ] - \lambda (x + y) + \frac {1}{2} \lambda^ {2} [ x, y ] & \\ + \frac {1}{2} [ \lambda (x + y), - \lambda (x + y) ] = \lambda^ {2} [ x, y ] \end{array}
\]

这证明了 (4)。

命题 5. 设 G 是李群，k 是 K 的非离散闭子域，G' 是将 G 看作定义在 k 上的李群所得的群，且 \(\phi\)（分别为 \(\phi'\)）是 G（分别为 G'）的指数映射。则 \(\phi\) 与 \(\phi'\) 在 0 的某个邻域上重合。

\(\phi\) 相对于 \(G'\) 满足定理4的假设 (i)，因而是 \(G'\) 的指数映射。

命题 6. 设 G 是李群芽，L 是其李代数，\(\phi: V \to G\) 是 G 的指数映射。对任意 \(x \in V\)，将 \(\mathrm{T}_x(L)\) 与 L 识别，从而 \(\varpi(x)\) 是 \(\phi\) 在 \(x\) 处的右微分，是从 L 到 L 的线性映射。当 \(x\) 充分接近 0 时，

\[
\varpi (x) = \sum_ {n \geqslant 0} \frac {1}{(n + 1) !} (\mathrm{ad} x) ^ {n}.
\]

给 \(\mathbf{L}\) 赋予一个与其拓扑相容且满足 \(\| [x,y]\| \leqslant \| x\| \| y\|\)（对所有 \(x,y\in L\)）的范数。只需考虑 \(G\) 是由 \(L\) 定义的李群芽且 \(\phi = \mathrm{Id}_{\mathrm{G}}\) 的情形。按定义，此时 \(\varpi (x)\) 是在 \(x\) 处的映射 \(y\mapsto y.x^{-1}\) 的切映射，该映射从 \(G\) 到 \(G\)。若 \(H(X,Y)\) 表示 Hausdorff 级数，则 \(\varpi (x)\) 在 \(\| x\|\) 充分小时，是在 0 处的映射 \(y\mapsto H(x + y, - x)\) 的切映射，该映射从 \(G\) 到 \(G\)。在 \(H(X + Y, - X)\) 中，关于 \(Y\) 的一次项之和为

\[
\sum_ {m \geqslant 0} \frac {1}{(m + 1) !} (\mathrm{adX}) ^ {m} \mathrm{Y}
\]

（第 II 章 § 6，第5号，命题5）。命题于是由《可微与解析流形》R，3.2.4 和 4.2.3 推出。

设 G 是李群芽且 \(t \in K\)。G 的 t 次幂映射是这样一个映射：它在 e 的某个开邻域上定义且解析，取值于 G，并在 e 的某个邻域上与映射

\[
g \mapsto \phi (t \phi^ {- 1} (g))
\]

重合，其中 \(\phi\) 是 G 的单射指数映射。

命题 7. (i) 若 \(t \in \mathbf{Z}\)，则 \(t\) 次幂映射在 \(e\) 的某个邻域上与映射 \(g \mapsto g^t\) 重合。

\begin{enumerate}
\def\labelenumi{(\roman{enumi})}
\setcounter{enumi}{1}
\item
  在 \(e\) 处，\(t\) 次幂映射的切映射是比值为 \(t\) 的伸缩。
\item
  若 \(h\) 是 \(t\) 次幂映射，\(h'\) 是 \(t'\) 次幂映射，定义在 \(G\) 上，则 \(h \circ h'\) 是 \((tt')\) 次幂映射，而 \(g \mapsto h(g)h'(g)\) 是 \((t + t')\) 次幂映射。
\item
  若 \(h\) 是 \(t\) 次幂映射且 \(u \in \mathbf{U}^n(\mathbf{G})\)，则 \(h_*(u) = t^n u\)。
\end{enumerate}

只需在 G 是由完备赋范李代数定义的李群芽，且所考虑的 t 次幂映射使用指数映射 \(\phi = \mathrm{Id}_{\mathrm{G}}\) 构造时证明命题。但在这种情形下一切都是显然的。

\subsection{指数映射的函子性}

命题 8. 设 G 和 H 是李群芽，h 是从 G 到 H 的态射，\(\phi_{\mathrm{G}}\) 和 \(\phi_{\mathrm{H}}\) 分别是 G 和 H 的指数映射。存在 L(G) 中 0 的一个邻域 V，使得 \(h \circ \phi_{\mathrm{G}}\) 与 \(\phi_{\mathrm{H}} \circ \mathrm{L}(h)\) 在 V 上重合。

给 \(\mathrm{L}(\mathrm{G})\) 和 \(\mathrm{L}(\mathrm{H})\) 赋予分别定义其拓扑且满足 \(\| [x,y]\| \leqslant \| x\| \| y\|\) 对所有 \(x\) 和 \(y\) 成立的范数。可以假设 \(\mathrm{G}\)（分别为 H）是由 \(\mathrm{L}(\mathrm{G})\)（分别为 \(\mathrm{L}(\mathrm{H})\)）定义的李群芽，从而 \(\phi_{\mathrm{G}}\)（分别为 \(\phi_{\mathrm{H}}\)）在 0 的某个邻域上与 \(\mathrm{Id}_{\mathrm{G}}\)（分别为 \(\mathrm{Id}_{\mathrm{H}}\)）重合。另一方面，\(\mathrm{L}(\mathrm{G})\) 中存在 0 的一个对称开邻域 W，使得 \(\mathrm{L}(h)\) 是从李群芽 W 到 H 的态射。根据定理1，\(\mathrm{L}(h)\) 在 0 的某个邻域上与 \(h\) 重合，命题得证。

粗略地说，若借助指数映射在单位元的某个邻域上将 G 和 H 分别与 L(G) 和 L(H) 识别，则从 G 到 H 的每个态射在 0 的某个邻域上都是线性的。

推论 1. 设 G 是李群芽，G' 是 G 的李子群芽，\(\phi\) 是 G 的指数映射。

\begin{enumerate}
\def\labelenumi{(\roman{enumi})}
\item
  存在一个开邻域 \(\mathbf{V}\)，它位于 \(\mathrm{L}(\mathbf{G}')\) 中的 0 附近，使得 \(\phi | \mathbf{V}\) 是从流形 \(\mathbf{V}\) 到 \(e\) 在 \(\mathbf{G}'\) 中的一个开邻域的同构。
\item
  设 \(x \in \mathrm{L}(\mathrm{G})\)。以下条件等价：(a) \(x \in \mathrm{L}(\mathrm{G}')\)；(b) \(\phi(\lambda x) \in \mathrm{G}'\)，当 \(|\lambda|\) 充分小时。
\item
  将命题8应用于从 \(\mathbf{G}'\) 到 \(\mathbf{G}\) 的典范嵌入即可得到，且 (ii) 由 (i) 推出。
\end{enumerate}

推论 2. 设 G 是李群，\(\rho\) 是 G 的解析线性表示，\(\phi\) 是 G 的指数映射。存在 L(G) 中 0 的一个邻域 V，使得

对所有 \(x \in V\) 均成立。

\[
\rho (\phi (x)) = \exp (\mathrm{L} (\rho) x)
\]

这由命题8及第3号的例2推出。

推论 3. 设 G 是李群，\(\phi\) 是 G 的指数映射。

\begin{enumerate}
\def\labelenumi{(\roman{enumi})}
\tightlist
\item
  存在 L(G) 中 0 的一个邻域 V，使得
\end{enumerate}

\[
\operatorname{Ad} (\phi (x)) = \exp \operatorname{ad} x.
\]

对所有 \(x \in V\) 均成立。

\begin{enumerate}
\def\labelenumi{(\roman{enumi})}
\setcounter{enumi}{1}
\tightlist
\item
  若 \(g \in \mathbf{G}\)，则存在 0 的一个邻域 \(\mathbf{W}\)，它位于 \(\mathrm{L}(\mathbf{G})\) 中，使得
\end{enumerate}

\[
g \phi (x) g ^ {- 1} = (\mathrm{Ad} g. x)
\]

对所有 \(x \in W\) 均成立。

\begin{enumerate}
\def\labelenumi{(\roman{enumi})}
\tightlist
\item
  这由推论2及 § 3 第12号命题44推出。
\end{enumerate}

\[
g.
\]

\subsection{子群上的诱导结构}

引理 4. 设 G 是有限维李群，\(\Omega\) 是 G 中 \(e\) 的对称开邻域，H 是 \(\Omega\) 的一个包含 \(e\) 的子集，并满足：条件 \(x \in \mathrm{H}\)、\(y \in \mathrm{H}\)、\(xy^{-1} \in \Omega\) 推出 \(xy^{-1} \in \mathrm{H}\)。设 \(r \in \mathrm{N}_{\mathbb{K}}\)。对所有 \(x \in \mathrm{H}\)，令 \(\mathfrak{h}_x\) 为所有满足下述性质的 \(a \in \mathrm{T}_x(\mathrm{G})\) 的集合：存在 K 中 0 的一个开邻域 I 及一个映射 \(f\)，它是从 I 到 G 的 \(\mathrm{C}^r\) 类映射，使得 \(f(0) = x, f(I) \subset H\)，且 \((\mathrm{T}_0f)(1) = a\)。

\begin{enumerate}
\def\labelenumi{(\roman{enumi})}
\item
  令 \(\mathfrak{h}_e = \mathfrak{h}\)。则 \(\mathfrak{h}\) 是 \(\mathrm{L}(\mathrm{G})\) 的李子代数，并且在 \(\mathrm{Ad}_{\mathrm{L}(\mathrm{G})}(\mathrm{H})\) 下不变。
\item
  \(\mathfrak{h}_x = x\mathfrak{h} = \mathfrak{h}x\) 对所有 \(x\in \mathrm{H}\) 成立（其中 \(x\mathfrak{h}\) 和 \(\mathfrak{h}x\) 在 \(\mathrm{T(G)}\) 中计算）。
\item
  设 \(\mathbf{V}\) 是 \(\mathbf{C}^r\) 类流形，\(v_0\) 是 \(\mathbf{V}\) 中的一点，且有一个映射 \(f\)，它是 \(\mathbf{C}^r\) 类的、从 \(\mathbf{V}\) 到 \(\mathbf{G}\) 的映射，满足 \(f(v_0) = e\) 且 \(f(\mathbf{V}) \subset \mathbf{H}\)。对于每个李子群芽 \(\mathbf{H}'\)，它属于 \(\mathbf{G}\)，且满足 \(\mathfrak{h}, f(v) \in \mathbf{H}'\)；当 \(v\) 充分接近 \(v_0\) 时这一关系成立。
\item
  对于每个李子群芽 \(\mathbf{H}'\)，它是 \(\mathbf{G}\) 的且李代数为 \(\mathfrak{h}\)，\(\mathbf{H}' \cap \mathbf{H}\) 是 \(e\) 在 \(H'\) 中的一个邻域。
\item
  对所有 \(x \in \mathbf{H}\) 及所有 \(a \in \mathfrak{b}_x\)，存在 K 中 0 的一个开邻域 I 及一个从 I 到 G 的映射 \(f\)，它属于 \(\mathbf{C}^{\omega}\) 类，使得 \(f(0) = x, f(I) \subset H, (T_0f)(1) = a\)。显然 \(\mathrm{K}\mathfrak{b} = \mathfrak{b}\)，并且 \(x\mathfrak{b}_y z = \mathfrak{b}_{xyz}\) 在 \(x, y, xy, xyz\) 属于 H 时成立。这推出 (ii)，并推出 \(\mathfrak{b}\) 在 \(\mathrm{Ad}_{\mathrm{L(G)}}(\mathrm{H})\) 下不变。
\end{enumerate}

设 \(a_1, a_2\) 属于 \(\mathfrak{h}\)。设 I 是 K 中 0 的一个开邻域，\(f_1, f_2\) 是从 I 到 G 的 \(\mathrm{C}^r\) 类映射，满足 \(f_j(0) = e, f_j(I) \subset H\)，\((\mathrm{T}_0 f_j)(1) = a_j (j = 1, 2)\)。定义 \(f: I \to G\) 为 \(f(\lambda) = f_1(\lambda)f_2(\lambda)\)。则 \(f\) 属于 \(\mathrm{C}^r\) 类且 \(f(0) = e\)。必要时缩小 I，可使 \(f(I) \subset H\)。另一方面，从 \(\mathrm{T}_e(G) \times \mathrm{T}_e(G)\) 到 \(\mathrm{T}_e(G)\) 的映射 \((g, g') \to gg'\) 在单位元处的切映射是加法；因此 \((\mathrm{T}_0 f) = a_1 + a_2\)。所以 \(a_1 + a_2 \in \mathfrak{h}\)，且 \(\mathfrak{h}\) 是 L(G) 的向量子空间。由于 \(x\mathfrak{h}x^{-1} = \mathfrak{h}\) 对所有 \(x \in H\) 都成立，\((\mathrm{Ad}f_1(\lambda)).a_2 \in \mathfrak{h}\) 对所有 \(\lambda \in I\) 都成立。根据 § 3 第12号命题44，映射 \(\lambda \mapsto \mathrm{Ad}f_1(\lambda)\) 在 0 处的切映射是映射 \(\lambda \mapsto \mathrm{ad}(\lambda a_1)\)；于是

\[
\left[ a _ {1}, a _ {2} \right] = (\operatorname{ad} a _ {1}). a _ {2} \in \mathfrak {h}
\]

由于 \(\mathfrak{h}\) 在 \(\mathrm{L}(\mathrm{G})\) 中闭，所以我们证明了 (i)。在证明的其余部分，固定一个李子群芽 \(\mathrm{H}'\)，它属于 \(\mathrm{G}\) 且李代数为 \(\mathfrak{h}\)。

设 \(V, v_0, f\) 如 (iii) 中。设 \(Y\) 是 \(G\) 中与 \(\mathfrak{h}\) 相关的左叶状结构（第1号）。对所有 \(y \in H'\)，\(T_y(H') = y\mathfrak{h}\)。另一方面，对所有 \(v \in V\)，\(T_v(V)\) 在 \(T_v(f)\) 下的像包含于 \(\mathfrak{h}_{f(v)} = f(v)\mathfrak{h}\)（按 \(\mathfrak{h}_{f(v)}\) 的定义）。根据《可微与解析流形》R，9.3.2，\(f\) 是从 \(V\) 到 \(Y\) 的态射。由于 \(H'\) 是 \(Y\) 的一片叶（《可微与解析流形》R，9.2.8），\(f(v) \in H'\) 当 \(v\) 充分接近 \(v_0\) 时。

设 \((a_{1},\ldots ,a_{s})\) 是 \(\mathfrak{h}\) 的一组基。存在 K 中 0 的一个开邻域 I 及映射 \(f_{1},\ldots ,f_{s}\)，它们是从 I 到 G 的 \(\mathbf{C}^r\) 类映射，满足 \(f_{j}(0) = e\)、\(f_{j}(\mathrm{I})\subset \mathrm{H}\)、\((\mathrm{Tf}_j)1 = a_j\)，对所有 \(j\)。根据 (iii)，\(f_{j}(\lambda)\in \mathrm{H}'\) 当 \(|\lambda |\) 充分小时。因此，\(f_{1}(\lambda_{1})f_{2}(\lambda_{2})\ldots f_{s}(\lambda_{s})\) 在 \(|\lambda_1|,\left|\lambda_2\right|,\ldots ,|\lambda_s|\) 充分小时构成 \(e\) 在 \(\mathbf{H}'\) 中的一个邻域；且这个邻域包含于 \(\mathbf{H}\)。所以得到 (iv)。

若 \(a \in \mathfrak{h}\)，则存在 0 的一个开邻域 \(\mathbf{I}\)，它位于 \(\mathbf{K}\) 中；以及一个映射 \(f\)，它是 \(\mathbf{C}^{\omega}\) 类的，从 \(\mathbf{I}\) 到 \(\mathbf{G}\)，使得 \(f(0) = e, f(\mathbf{I}) \subset \mathbf{H}'\)，\((\mathrm{T}_0f)1 = a\)。结合 (iv)，这推出 (v)。

定义 2. 称 \(\mathfrak{h}\) 为 H 在 \(e\) 处的切子代数。

命题 9. 设 G 是有限维李群，H 是 G 的子群。

\begin{enumerate}
\def\labelenumi{(\roman{enumi})}
\item
  H 上存在唯一的解析流形结构，具有如下性质：对 1 与 \(\omega\) 之间的任意 r、任意 \(C^r\) 类流形 V 及任意从 V 到 H 的映射 f，\(f\) 作为从 V 到 H 的映射属于 \(C^r\) 类，当且仅当 f 作为从 V 到 G 的映射属于 \(C^r\) 类。
\item
  在此结构下，\(\mathbf{H}\) 是李群，其典范嵌入 \(i\) 把 \(\mathbf{H}\) 映入 \(\mathbf{G}\)，是浸入，并且 \(\mathbf{L}(i)(\mathbf{L}(\mathbf{H}))\) 是 \(e\) 处 \(\mathbf{H}\) 的切李子代数。
\end{enumerate}

在 (i) 中，唯一性显然。我们证明存在性。设 \(\mathfrak{b}\) 是 H 在 \(e\) 处的切李代数。设 H' 是 G 的李子群芽，其李代数为 \(\mathfrak{b}\)。以 H' 的一个开子群芽替换 H' 后，可假设 H' ⊆ H（引理4 (iv)）。对所有 \(x \in \mathrm{H}\)，\(xH'x^{-1}\) 是 G 的李子群芽，其李代数为 \(x \mathfrak{h}x^{-1} = \mathfrak{b}\)。因此 H' ∩ \((xH'x^{-1})\) 在 H' 中开（第2号定理3），且映射 y → \(xyx^{-1}\) 是从 H' ∩ \(x^{-1}H'x\) 到 \(xH'x^{-1}\) ∩ H' 的同构。利用 § 1 第9号命题18，存在 H' 的一个开李子群芽 W 以及 H 上的一个李群结构，具有如下性质：W 在 H 中开，并且 H 与 H' 上的流形结构在 W 上诱导出相同的结构。于是，H 到 G 的典范嵌入 i 是浸入，且 L(i)(L(H)) = L(H') = \(\mathfrak{b}\)。此外，设 V 和 f 如 (i) 中。若 f: V → H 属于 \(C^r\) 类，则 \(i \circ f\): V → G 属于 \(C^r\) 类。反之，假设 \(i \circ f\): V → G 属于 \(C^r\) 类；我们证明 f: V → H 属于 \(C^r\) 类。通过平移，只需考虑存在 v0 ∈ V 使得 f(v0) = e，并证明 f: V → H 在 v0 的某个开邻域上属于 \(C^r\) 类的情形。现在，根据引理4 (iii)，当 v 充分接近 v0 时，f(v) ∈ H'，所述断言随即成立。因此我们证明了 (i)，而 (ii) 也在此过程中得到。

定义 3. 命题9中定义在 H 上的李群结构称为由 G 上的李群结构在 H 上诱导的结构。

若 H 是 G 的李子群，则 H 的李群结构由 G 上的李群结构诱导（《可微与解析流形》，R，5.8.5）。

若 \(\mathbf{G} = \mathbf{R}\) 且 \(\mathbf{H} = \mathbf{Q}\)，则 \(\mathfrak{h} = \{0\}\)，因而 \(\mathbf{H}\) 上的诱导结构是离散李群结构。同样，若 \(\mathbf{G} = \mathbf{C}\)（看作复李群）且 \(\mathbf{H} = \mathbf{R}\)，结论也成立。
